% ==========================================================================
%  Context Selection Is a Partial Order
%  Assembled arXiv article (editor build).
%
%  AUTHOR NOTE ON CITATIONS (see also the note preceding the bibliography and
%  the .bib header): only Blackwell (1953), Le Cam (1964/1986), and
%  Ethayarajh, Choi & Swayamdipta (2022) are author-verified. EVERY other
%  reference was surfaced by automated literature search and is marked
%  [TO VERIFY]; the bibliography MUST be checked against real arXiv/DBLP
%  before submission.
% ==========================================================================
\documentclass[11pt]{article}
\usepackage[T1]{fontenc}     % scalable Type1 font encoding
\usepackage{lmodern}         % scalable Latin Modern fonts (required for microtype expansion)
\usepackage[margin=1in]{geometry}
\usepackage{amsmath,amssymb,amsthm,booktabs,hyperref,natbib,microtype,enumitem}

\theoremstyle{plain}
\newtheorem{theorem}{Theorem}
\newtheorem{proposition}{Proposition}
\newtheorem{corollary}{Corollary}
\newtheorem{lemma}{Lemma}
\theoremstyle{definition}
\newtheorem{definition}{Definition}
\newtheorem{assumption}{Assumption}
\theoremstyle{remark}
\newtheorem{remark}{Remark}

%, convenience macros for consistent notation, 
\newcommand{\PASS}{\mathrm{PASS}}
\newcommand{\dhat}{\widehat{\delta}}
\newcommand{\Dcal}{\mathcal{D}}
\newcommand{\Wcal}{\mathcal{W}}
\newcommand{\Scal}{\mathcal{S}}
\newcommand{\Fcal}{\mathcal{F}}
\newcommand{\Vcal}{\mathcal{V}}
\newcommand{\Bord}{\succeq_{B}}

\title{Context Selection Is a Partial Order:\\
A Blackwell Theory of Why No Relevance Score Transfers Across Tasks}

\author{%
  Andriy Bilous \\
  Department of Information Systems and Networks \\
  Lviv Polytechnic National University, Lviv, Ukraine \\
  ORCID: 0009-0006-5467-7932 \\
  \texttt{andriy.bilous@uitware.com}
}

\date{}

\begin{document}
\maketitle

% ==========================================================================
\begin{abstract}
% ==========================================================================
Choosing what to put in a language model's prompt (a retrieved passage, a
private convention, a repository file, a demonstration) is almost universally
posed as a \emph{scalar} ranking problem: score each candidate by relevance,
mutual information, $\Vcal$-usable information, or an information-gain
proxy, and take the top-$k$. We argue this primitive is structurally wrong.
Modeling each context source as a statistical \emph{experiment} on the task's
latent requirement, we transport Blackwell's comparison of
experiments~\citep{blackwell1953} and Le~Cam's deficiency~\citep{lecam1964,lecam1986}
to context selection. The resulting order over sources is \emph{partial}, and
our central claim, an application-level consequence of Blackwell's theorem, is
an \emph{incomparability} result: when two sources are Blackwell-incomparable
(neither is a garbling of the other), no scalar score can rank them correctly
across all tasks. This gives a formal account of the field's
``relevance-does-not-transfer'' folklore. To make the order testable on a real
model with a real verifier, we define a decision-rule-restricted deficiency
\emph{surrogate}, a value gap evaluated at the fixed decision rule the model
induces, and equip it with a distribution-free \emph{uniform certificate}
(exact Clopper--Pearson intervals with a Bonferroni correction over the task
battery), which we show is the multiplicity correction a naive
Dvoretzky--Kiefer--Wolfowitz bound omits. We instantiate the theory on six
language models across four vendors (Anthropic, OpenAI, DeepSeek, Moonshot).
On all six we certify ($\geq 90\%$ confidence, both directions) that two private
context sources are empirically incomparable, and a query-conditioned relevance
score mis-ranks them on a confound-controlled trap task. We further report two
findings the framework anticipates and the surrogate detects: \textbf{(i)
anti-monotonicity}, supplying a strict superset of a sufficient source can
sharply degrade performance, a certified collapse from $100\%$ to as low as
$0\%$ pass rate, which falsifies the monotonicity that scalar context-value
scores silently assume; and \textbf{(ii)} a \textbf{model-relative value of
private context}, $I(S;W\mid\theta)$, that varies continuously across models
(the guessability of a fixed convention spans $0\%$--$75\%$ across vendors) and
tracks a vendor's training rather than raw capability. The anti-monotonicity effect is large and robust across the
capability ladder but is \emph{not} monotone in capability, and the empirical
program rests on a single adversarially constructed source pair under a
restricted estimator. The conclusion is methodological: context
selection should be treated as a partial-order problem, and a certified
deficiency surrogate, not a scalar, is the right operational object. We give
the selection rule this implies, \textsc{Probe-Select}, which keeps the certified
non-dominated set and diversifies over incomparable coordinates; it provably never
retains a dominated source, recovers top-$k$ when sources are totally ordered, and
on our measured data avoids the anti-monotonicity trap a superset heuristic falls
into. Scaling it beyond curation-size candidate sets to corpus retrieval is the
central problem we leave open.
\end{abstract}

% ==========================================================================
\section{Introduction}
\label{sec:intro}
% ==========================================================================

Production language-model systems spend much of their budget deciding
\emph{what context to place in the prompt}. Retrieval-augmented generation ranks
passages, in-context learning selects demonstrations, and agentic coding tools
choose which files and conventions to surface. Across all of these, the
operative tool is a \textbf{scalar score}, lexical or embedding relevance,
mutual information, $\Vcal$-usable information~\citep{ethayarajh2022},
expected information gain, used to impose a total order and retain the top
candidates. The scalar is convenient: it sorts, thresholds, and composes. We
argue it is also the wrong abstraction, and that a classical apparatus from
statistical decision theory says precisely why.

\paragraph{A folklore observation, restated as structure.}
A persistent, informal observation undercuts the scalar paradigm: \emph{relevance
does not transfer across tasks}. A source that helps one task hurts another;
``the best context'' is task-dependent. This is usually treated as a modeling
nuisance to be engineered around with better rankers. We show it is the
signature of a \emph{partial order}. We model a context source $W$ as a
statistical experiment on the task's latent requirement $S$ (the object a
verifier checks; \S\ref{sec:formalization}), which lets us import two classical
tools. Blackwell's order ranks experiments by garbling: $W_1 \Bord W_2$ iff
$W_2$'s signal is a noised (garbled) version of $W_1$'s, and, by the
Blackwell--Sherman--Stein theorem, this holds \emph{iff} an idealized
decision-maker does at least as well with $W_1$ as with $W_2$ on
\emph{every} decision problem~\citep{blackwell1953,torgersen1991}. The order is
\emph{partial}: many source pairs are incomparable, neither garbling the other.
Incomparability is exactly the regime in which no scalar summary can be faithful,
because a scalar is a total order and a total order extending a strict partial
order must invert at least one incomparable pair on some task.

\paragraph{What is and is not new.}
The mathematical machinery, the Blackwell garbling order, the
Blackwell--Sherman--Stein equivalence, Le~Cam's deficiency, is classical and we
claim none of it as new. Our contribution is an \emph{application}: transporting
comparison-of-experiments to LLM context selection; an incomparability
\emph{framing} that turns cross-task non-transfer of relevance from folklore into
a stated consequence; and an \emph{estimator}, a certified, decision-restricted
deficiency surrogate, that makes the partial order measurable on a real model
behind a real verifier. One gap the framing must
respect, and that we close, concerns the decision rule. Blackwell's theorem is a statement
about the \emph{optimal} (Bayes) use of an experiment, whereas every quantity we
measure, the pass rate
$\PASS(W,T) = \mathbb{E}\!\left[v_T(\delta_M(T,w))\right]$, is evaluated
at a \emph{single, fixed, sub-optimal} decision rule $\delta_M$ induced by a
particular model $M$. ``The optimal user of $W_2$ beats the optimal user of
$W_1$ on some task'' does not by itself imply ``the fixed model does better with
$W_2$ than with $W_1$.'' We therefore develop the order and the incomparability
statement in a $\delta_M$-restricted category (\S\ref{sec:formalization},
\S\ref{sec:incomparability}), prove a one-way bridge to the idealized Blackwell
order, and confine every empirical claim to the restricted object that the
experiments actually probe. Likewise, our estimator
$\dhat_{\Dcal}$ performs no minimization over garblings, so it is a
\emph{decision-restricted value-deficiency surrogate}, not a Le~Cam deficiency;
we name it as such and relate the two formally (\S\ref{sec:estimator}) rather
than borrowing the deficiency's pedigree.

\paragraph{A latent-coordinate model that makes the order non-vacuous.}
A theory referee will object, correctly, that our sources are \emph{fixed
deterministic strings} that do not vary with the task instance, so a kernel
$\kappa_W(\cdot\mid S)$ would be a constant, hence a Blackwell-trivial
(uninformative) experiment with a trivial $\sigma$-algebra. We resolve this by
modeling the latent requirement as a \emph{vector of coordinates}
$S = (S_1, S_2, \dots)$ and each source as a deterministic experiment observing a
\emph{projection} $\kappa_{W_i}(\cdot\mid S) = \delta_{\pi_i(S)}$. Comparison of
experiments then reduces to refinement of the projections' sub-$\sigma$-algebras
(Blackwell for deterministic experiments), and incomparability becomes the
statement that neither projection refines the other, a property we establish for
our concrete pair (\S\ref{sec:formalization}, \S\ref{sec:experiments}). Under
this construction the garbling argument is rigorous rather than a metaphor about
the logical content of a string.

\paragraph{Empirical findings.}
We instantiate the theory on two private, equally on-topic conventions about the
same module: an API call contract $W_1$ and a wire-encoding invariant $W_2$
(\S\ref{sec:experiments}), with a four-cell verifier battery using executable
PASS$/$FAIL judgments (return code $0$ is PASS). This imposes a
restriction: the battery is a fixed, adversarially constructed
opposition, not a sampled task distribution, and the certified incomparability is
carried by one API cell and one encoding cell; we report the effective support
explicitly and do not call $\Dcal$ a ``distribution.'' On this instrument
we obtain, on \emph{real} model calls with no modeled numbers:

\begin{itemize}
  \item \textbf{Certified empirical incomparability on all six models.} The
  uniform certificate clears $\geq 90\%$ confidence in \emph{both} directions:
  the two-sided lower bounds $L_{\Dcal}(W_1,W_2)\,/\,L_{\Dcal}(W_2,W_1)$
  are
  Haiku-4.5 $+0.88\,/\,+0.83$ ($n{=}80$),
  Sonnet-4.6 $+0.76\,/\,+0.76$ ($n{=}40$),
  Opus-4.8 $+0.55\,/\,+0.55$ ($n{=}20$),
  GPT-5.5 $+0.76\,/\,+0.76$ ($n{=}40$),
  DeepSeek-V4-Pro $+0.347\,/\,+0.715$ ($n{=}40$), and
  Kimi-K2.6 $+0.675\,/\,+0.762$ ($n{=}40$). Incomparability is therefore not a
  single-lab artifact.
  \item \textbf{A real relevance score mis-ranks the pair.} On a trap task
  saturated with encoding vocabulary, query-conditioned lexical relevance scores
  $\varphi(W_1){=}0.36 < \varphi(W_2){=}0.44$ (ranking $W_2 \succ W_1$), yet
  realized PASS ranks $W_1 \gg W_2$ (e.g.\ $96$--$100\%$ vs.\ $0\%$). On the
  honest tasks relevance agrees with PASS; it fails \emph{only} where utility
  diverges from topicality, the precise prediction of the incomparability
  framing. We state plainly that this is shown for bag-of-words lexical
  relevance, the scalar the simplest practitioners use, and is not yet replicated
  with a learned reranker.
  \item \textbf{Certified anti-monotonicity.} The superset $W_{1+} = W_2 +\!\!+\, W_1$,
  which literally \emph{contains} $W_1$'s signal, sharply degrades the API cells:
  Sonnet and Kimi collapse $100\% \to 0\%$ (interference $\Psi_T = -1.00$,
  certified upper bound $\leq -0.69$); GPT-5.5 certifies $\Psi_T \leq -0.60$ on
  \texttt{api\_post\_ok}. The interference certificate clears on at least one
  cell for five of six models (Opus reproduces, $\Psi_T = -0.75/-0.80$, but
  $n{=}20$ is too few to certify; we give the power argument in
  \S\ref{sec:experiments}). An order control shows the harm does not recover
  under reordering, so it is not a lost-in-the-middle artifact. One
  caveat: the effect is large and robust across the capability ladder but is
  \emph{not} monotone in capability (Sonnet, the middle model, is the starkest;
  Opus, the strongest, less so), so we claim no capability trend.
  \item \textbf{A model-relative value of private context.} The \texttt{none}-arm
  baseline on \texttt{api\_argorder} forms a continuous spectrum, Claude
  $\sim\!0\%$, DeepSeek $2\%$, Opus $15\%$, Kimi $45\%$, GPT-5.5 $75\%$, so the
  value $I(S;W\mid\theta)$ of the private context (where $\theta$ denotes the
  model together with its pretraining prior over conventions) varies
  \emph{continuously} across models and tracks a vendor's training rather than
  raw capability (DeepSeek is tight like Claude despite a different vendor). This
  is an empirical instance of the theory's $\theta$-dependence, and it caps which
  cells can be certified for which model.
\end{itemize}

\paragraph{Methodological asymmetry.}
The Anthropic models run agentically (\texttt{claude -p}, up to six turns); the
non-Anthropic models run single-shot (completion plus code extraction via an
OpenAI-style API). We flag this throughout: vendor is confounded with harness, so
every \emph{cross-vendor magnitude} comparison (the value spectrum, the
interference ladder) inherits that confound, whereas the binary within-model
claims do not. The single-shot setting is arguably a \emph{cleaner} probe of
whether a model uses supplied context, since it removes the agentic-exploration
confound; we treat the asymmetry as a limitation to be reported, not explained
away.

\paragraph{Contributions.}
\begin{enumerate}
  \item[\textbf{(C1)}] \textbf{A decision-independent partial order over context
  sources} (\S\ref{sec:order}). We instantiate the Blackwell garbling order on
  context sources and state the for-all-tasks dominance guarantee for the
  idealized decision-maker via the Blackwell--Sherman--Stein theorem, together
  with the $\delta_M$-restricted order the experiments probe and a one-way
  bridging lemma between them. This is strictly stronger than any scalar
  relevance score, which is task-specific.
  \item[\textbf{(C2)}] \textbf{An incomparability result for context selection}
  (\S\ref{sec:incomparability}). We give an application-level consequence of
  Blackwell's theorem: Blackwell-incomparable sources cannot be ranked by
  \emph{any} query-independent scalar $\varphi(W)$ for all tasks, and we
  separately establish, as an empirically grounded statement, that even a
  query-\emph{conditioned} relevance score $\varphi(W,T)$ mis-ranks them on a
  constructed trap. The theorem covers source-level scalars,
  and the trap covers the query-conditioned case; the two are distinct. This
  is the first formal account, to our knowledge as of mid-2026, of cross-task
  non-transfer of relevance.
  \item[\textbf{(C3)}] \textbf{A measurable, certified deficiency surrogate}
  (\S\ref{sec:estimator}). We define
  $\dhat_{\Dcal}(W_1,W_2) = \max\!\big(0,\,
  \sup_{T\in\Dcal}[\PASS(W_2,T) - \PASS(W_1,T)]\big)$, a
  decision-restricted value-deficiency surrogate (named as such, not as a Le~Cam
  deficiency), and a distribution-free uniform certificate: exact
  Clopper--Pearson intervals~\citep{clopperpearson1934} at Bonferroni
  level~\citep{bonferroni1936} $\alpha = \eta/(2|\Dcal|)$, valid jointly on
  a single event, so that $L_{\Dcal}>0$ certifies a positive value gap at
  confidence $\geq 1-\eta$. We argue why a naive DKW
  bound~\citep{dvoretzky1956} under-covers a supremum over $|\Dcal|$ cells,
  and we give the matching interference certificate that the anti-monotonicity
  finding requires.
  \item[\textbf{(C4)}] \textbf{A constructive partial-order selector}
  (\S\ref{sec:method}). \textsc{Probe-Select} keeps the certified non-dominated set and then
  diversifies by certified task coverage. We prove it (i) never retains a Blackwell-dominated source
  (monotone safety), (ii) recovers scalar top-$k$ when the sources are totally ordered, and (iii)
  otherwise covers distinct latent coordinates with the standard $(1-1/e)$ submodular guarantee. On
  our measured data it returns the $\{W_1, W_{1+}\}$ front and rejects the superset-only choice that
  the anti-monotonicity result shows would collapse, answering ``what replaces the scalar,'' with
  an explicit account of its curation-scale cost and the corpus-scale open problem.
  \item[\textbf{(C5)}] \textbf{A cross-vendor empirical program} (\S\ref{sec:experiments}).
  Certified incomparability and the relevance mis-rank on six models across four
  vendors; certified anti-monotonicity on five of six; and a continuous map of
  the model-relative value of private context, all on real model calls, with the
  scope limits (single source pair, vendor--harness confound,
  capability-non-monotonicity) stated in \S\ref{sec:limitations}.
\end{enumerate}

\paragraph{Roadmap.}
\S\ref{sec:related} positions the work against scalar context-selection criteria
and the recent transport of comparison-of-experiments into machine learning, and
frames our delta as application-plus-framing-plus-estimator rather than new core
mathematics. \S\ref{sec:formalization} gives the latent-coordinate experiment
model, defines $\theta$ and the $\delta_M$-restricted order, and states the
non-refinement (hence incomparability) construction. \S\ref{sec:order} develops
the partial order (C1) and the bridging lemma. \S\ref{sec:incomparability} proves
the source-level incomparability statement and isolates the query-conditioned
case (C2). \S\ref{sec:estimator} defines the surrogate, the uniform certificate,
and the interference certificate, and relates the surrogate to the true Le~Cam
deficiency (C3). \S\ref{sec:method} gives the \textsc{Probe-Select} selector and its
monotone-safety, top-$k$-generalization, and coverage guarantees, with a worked illustration on the
measured data (C4). \S\ref{sec:experiments} reports the six-model program with full
per-cell tables (C5). \S\ref{sec:limitations} states threats to validity and
scope, and \S\ref{sec:conclusion} concludes. Proofs accompany the statements
in-line; the bibliography is provisional and every non-classical reference is
marked \textsc{[to verify]} pending a check against arXiv/DBLP (see the author
note preceding the references).

% ==========================================================================
\section{Related Work}
\label{sec:related}
% ==========================================================================

\paragraph{Bibliography provenance (author note).}
Only three references in this paper have been verified by the authors against their
primary sources: Blackwell's comparison of experiments~\citep{blackwell1953},
Le~Cam's deficiency and the comparison of statistical experiments~\citep{lecam1964,lecam1986,torgersen1991},
and $\Vcal$-usable information~\citep{ethayarajh2022}. Every other
reference cited below was surfaced by an automated literature search and is marked
\texttt{[TO VERIFY]}; the corresponding \texttt{.bib} entries carry a \texttt{\% TO VERIFY}
comment. Our novelty claims and our characterization of prior art \emph{must} be re-checked
against arXiv/DBLP before submission, including a fresh search for any post-2025 work porting
statistical deficiency to language-model context. We flag this explicitly because our
own due-diligence pipeline previously over-read a candidate prior-art paper as containing
results it does not contain (see the discussion of~\citep{akdemir2025} below); the lesson is
that machine-summarized prior art is unreliable and a manual full-text read is mandatory.

\paragraph{Positioning.}
The machinery we use, Blackwell's garbling order~\citep{blackwell1953} and Le~Cam's
deficiency~\citep{lecam1964,lecam1986}, is classical, and we claim no new core mathematics.
Our contribution is an \emph{application}, a \emph{framing}, and an \emph{estimator}:
(i) we treat each candidate context source for a language model as a statistical experiment
on the task's latent requirement and rank sources by the Blackwell order, with a
decision-independent, for-all-tasks dominance guarantee (\S\ref{sec:order});
(ii) we observe, as an application-level corollary of Blackwell's theorem, that
Blackwell-\emph{incomparable} sources cannot be ranked by any scalar score for all tasks
(\S\ref{sec:incomparability}), which gives the first formal account of the
empirical folklore that ``relevance does not transfer across tasks''; and
(iii) we supply a decision-restricted Le~Cam deficiency \emph{surrogate} with a
distribution-free finite-sample certificate that is auditable on a real model with a real
verifier (\S\ref{sec:estimator}). To our knowledge, and subject to the caveat above,
no prior work applies the comparison of experiments to selecting inference-time context for
language models.

\subsection{Scalar context and demonstration selection}
\label{sec:related-scalar}

Information-theoretic context and demonstration selection for language models is dominated by
\emph{scalar} criteria, each of which induces a \emph{total} order over candidate sources and
takes a top-$k$ prefix. The dominant family scores topical relevance, either lexically or with
learned rankers trained jointly with the generator~\citep{rankrag2024}. A second family scores
\emph{usable} information: $\Vcal$-usable (predictive-$\Vcal$) information measures
how much a fixed model family can extract from an input~\citep{ethayarajh2022}, and recent work
uses it to enhance contextual faithfulness in language models~\citep{cale2025}. A third family scores expected information
or information gain, mutual information and Bayesian experimental design for language-model
querying~\citep{bedllm2026}, and information-pursuit variants with conformal
stopping~\citep{conformalip2025}. A fourth family attributes value to demonstrations through
game-theoretic or influence-based valuations~\citep{demoshapley2024,inficl2024}. A fifth treats
the question as binary classification: a ``sufficient-context'' detector predicts whether the
supplied context is enough to answer~\citep{sufficientcontext2024}.

Every criterion in this family collapses the comparison of two sources to a comparison of two
real numbers, and is therefore a \emph{scalar} in the precise sense of
Theorem~\ref{thm:incomparable}: it imposes one total order on $\{W_1,W_2\}$. Our delta is
structural, not a better scalar. We show (\S\ref{sec:incomparability}) that when two
sources are Blackwell-incomparable, as our two private conventions are
(\S\ref{sec:experiments}), \emph{any} such scalar, including relevance, mutual
information, $\Vcal$-usable information, perplexity-gap, and expected information gain,
must mis-rank them on some task. These criteria are not competitors to be out-performed but
instances of the class our incomparability result rules out: each is Blackwell-monotone, so
each is a one-dimensional projection of a partial order it cannot recover.
We substantiate the monotonicity claim for the closest neighbor, $\Vcal$-usable
information, in Proposition~\ref{prop:vinfo-monotone} rather than asserting it; we make the
weaker conjecture that the same holds for the remaining members of the family.
Two points sharpen the separation from the strongest neighbors. First, against expected
information gain and Bayesian experimental design~\citep{bedllm2026,conformalip2025}, the
relevant axis is \emph{decision-dependence}: those criteria are scalar scores optimized for a
particular query or stopping rule, whereas our Blackwell guarantee
(Theorem~\ref{thm:blackwell}) is decision-\emph{independent}, it holds for every prior, every
bounded utility, and every task simultaneously. Second, against query-conditioned relevance
$\varphi(W,T)$, which is not a fixed total order and is therefore \emph{not} ruled out by
Theorem~\ref{thm:incomparable}, we make a separate, empirical claim: a relevance-trap task
on which lexical relevance ranks $\varphi(W_2)>\varphi(W_1)$ while realized PASS ranks
$W_1\gg W_2$ (\S\ref{sec:experiments}). We are careful not to present the trap as covered
by the theorem; the theorem governs query-independent scalars, and the trap is the empirical
patch for the query-conditioned case.

\subsection{Task-dependence of context utility}
\label{sec:related-taskdep}

A parallel literature establishes that the utility of context is task-dependent rather than an
intrinsic property of the source, but does so only empirically or informally: some work argues
directly that topical relevance is not the same as answer
utility~\citep{relevanceutility2025,sureRAG2026}, and the broader phenomenon that adding more
or longer context can \emph{degrade} performance has been documented from several
angles~\citep{contexthurts2025,longcontextdegrade2026}. A separate strand observes
Pareto-incomparability of prompts, but along competing \emph{metrics within a single
task}~\citep{paretoprompts2025}, which is orthogonal to our object: incomparable
\emph{sources across tasks}. The folklore that the right context is task-dependent, that
relevance does not transfer across tasks, is exactly what our contribution converts into a
theorem. Cross-task non-transfer is not a modeling nuisance to be engineered around but a
necessary consequence of the Blackwell order being partial: any scalar imposes a total order
that the true partial order must violate somewhere. The empirical degradation results are
likewise reframed: our anti-monotonicity finding (\S\ref{sec:experiments}) is not merely
that ``more context can hurt'', an observation already in the
literature~\citep{contexthurts2025,longcontextdegrade2026}, but that a strict
\emph{superset} of a sufficient source can hurt, on sources we have separately \emph{certified}
to be Blackwell-incomparable, with a per-task interference certificate; the algebraic object
$\Psi_T<0$ is one that no monotone-submodular context-value scalar can carry.

\subsection{The comparison of experiments in machine learning}
\label{sec:related-deficiency}

The Blackwell--Le~Cam apparatus has recently entered machine learning, but exclusively for
ranking \emph{tasks} or \emph{domains}, never inference-time context for language models.
Statistical deficiency has been used to estimate task inclusion at training time, reconstructing
a classic NLP pipeline ordering over \emph{tasks}~\citep{fosse2025}; this orders the objects a
model is trained on, not the context supplied at inference, and provides neither an
incomparability result nor a PASS-rate certificate. A transfer-learning line builds a
deficiency-induced \emph{directional} dominance with a finite-sample proxy~\citep{akdemir2025};
we read this work in full and confirm it contains no incomparability statement, no formal
partial order (only a directional distance), and no distribution-free certificate (it uses an
MMD proxy that its own remarks disclaim as a guarantee), and that its domains are vision,
genomics, and reinforcement-learning control with no language-model, prompt, or retrieval
content. We therefore cite it as the deficiency-in-ML lineage for our order
(\S\ref{sec:order}), and differentiate on three axes: object (context sources vs.\
domains), structure (incomparability vs.\ directional dominance only), and a
distribution-free finite-sample certificate (which the transfer-learning deficiency work does
not provide). Comparison-of-experiments ideas have also been connected to loss design and
learning from corrupted labels~\citep{vanrooyen2014}; that line concerns the geometry of losses,
not representation or context selection, and we cite it only as evidence that the apparatus has
a foothold in learning theory.

\subsection{The Blackwell order in artificial intelligence}
\label{sec:related-llmrep}

Blackwell's theorems have recently been surveyed in an artificial-intelligence
context~\citep{blackwellrep2026}, evidence that the comparison-of-experiments apparatus is
recognized as relevant to the field; but to our knowledge no prior work applies the Blackwell
order to a language model's \emph{context sources}. We order external context sources supplied to
the model, prove the incomparability consequence, and carry out a certified empirical
Blackwell-incomparability test across six models and four vendors. To our knowledge no prior
work states a Blackwell partial order over context sources, derives the no-scalar-ranks
consequence for language models, or supplies a decision-restricted deficiency estimator with a
distribution-free certificate, the three pieces this paper contributes.

\subsection{Mathematical foundations}
\label{sec:related-foundations}

We rely on the standard foundations of the comparison of statistical experiments. The
equivalence between the garbling order and uniform dominance of Bayes risk is the
Blackwell--Sherman--Stein theorem~\citep{blackwell1953,torgersen1991}; deficiency and the metric
structure of the experiment ordering are due to Le~Cam~\citep{lecam1964,lecam1986}, with
Torgersen's monograph as the canonical reference~\citep{torgersen1991}. The fact that the
Blackwell order is a non-lattice partial order, so that genuinely incomparable pairs exist and
no scalar can be order-preserving, is also classical and is documented in the same
comparison-of-experiments literature~\citep{torgersen1991}. Our incomparability theorem
(Theorem~\ref{thm:incomparable}) is, mathematically, a short corollary of Blackwell's
theorem applied to a constructed pair of decision problems; we state it as an application-level
result about language-model context, not as new structure in the order itself. The genuinely
new contributions are the framing of context selection as a partial-order problem, the
decision-restricted PASS-rate deficiency surrogate, and its distribution-free uniform
certificate, none of which the neighbors above provide.

% ==========================================================================
\section{Formalization}
\label{sec:formalization}
% ==========================================================================

We model the choice of \emph{what context to supply} a language model as a comparison of
statistical experiments on the task's latent requirement. This section fixes the objects and,
crucially, separates two decision rules that the informal literature (and an earlier draft of
this work) silently conflate: the \emph{Bayes-optimal} rule, against which the classical
Blackwell order is defined, and the \emph{fixed model} rule, which is the object every empirical
quantity in this paper actually measures. Keeping these distinct is what makes the bridge from
theorem to experiment rigorous.

\subsection{Tasks, latent requirements, and context sources}

A \emph{task} $T$ fixes a latent requirement $S_T$, the object a deterministic verifier
$v_T$ checks (e.g.\ the intended program behaviour). We take $S_T$ to be a \emph{vector of
coordinates}
\begin{equation}
\label{eq:coordinates}
S \;=\; (S_1, S_2, \dots, S_m) \;\in\; \Scal \;=\; \Scal_1 \times \cdots \times \Scal_m,
\end{equation}
where each coordinate $S_j$ encodes one logically independent requirement (in our instantiation,
$S_1$ is the private call-contract convention and $S_2$ is the private wire-encoding convention;
see \S\ref{sec:experiments}). Modelling $S$ as a product of coordinates, rather than an
opaque atom, is the technical device that makes the experiment model well posed for
fixed-string sources; we return to this in \S\ref{subsec:sources} and Lemma~\ref{lem:noref}.

\begin{definition}[Context source as an experiment]
\label{def:source}
A \emph{context source} $W$ is a Markov kernel $\kappa_W : \Scal \to \Delta(\Wcal)$,
i.e.\ a statistical experiment indexed by the latent parameter $S$, emitting a signal
$w \sim \kappa_W(\cdot \mid S)$ in a signal alphabet $\Wcal$ (the source's text). The
experiment is identified with its kernel; the sample space is $\Wcal$ and the observed
datum is the supplied text $w$.
\end{definition}

\begin{definition}[Decision rules and PASS]
\label{def:decision}
A \emph{decision rule} is a (possibly randomized) map $\delta : (T,w) \mapsto \widehat S$
producing a candidate solution. Utility is the binary verifier $u(\widehat S, S) = v_T(\widehat S)
\in \{0,1\}$. Two decision rules matter:
\begin{enumerate}
  \item[(i)] the \emph{model rule} $\delta_M$ realized by a fixed language model $M$ at a fixed
    decoding configuration $\theta$ (model weights and sampling parameters; see
    \S\ref{subsec:theta}). Its realized pass rate on task $T$ under source $W$ is
    \begin{equation}
    \label{eq:pass}
    \PASS_M(W,T) \;:=\; \mathbb E\big[\, v_T\!\big(\delta_M(T,w)\big)\,\big],
    \qquad w \sim \kappa_W(\cdot \mid S_T);
    \end{equation}
  \item[(ii)] the \emph{Bayes rule} $\delta^\star_{T}$, the (utility-optimal) decision rule for
    task $T$ given a prior over $S$, achieving the optimal value
    $V^\star(W,T) := \sup_{\delta} \mathbb E[v_T(\delta(T,w))]$.
\end{enumerate}
\end{definition}

The distinction in Definition~\ref{def:decision} is load-bearing. Blackwell's theorem
(\S\ref{sec:order}) is a statement about $V^\star$ (the optimal user of an experiment); every
number we report, $\PASS_M$, the deficiency estimator of \S\ref{sec:estimator}, the
interference statistic of \S\ref{subsec:interference}, and the entire six-model program of
\S\ref{sec:experiments}, is a statement about the single fixed rule $\delta_M$. We therefore
develop the order in two registers and prove an explicit bridge between them
(Lemma~\ref{lem:bridge}).

\subsection{Sources as deterministic projection experiments}
\label{subsec:sources}

In our instantiation each context source is a \emph{fixed string}: a snippet stating the
private call contract ($W_1$), or the private wire-encoding invariant ($W_2$). A naive reading
treats $\kappa_W(\cdot\mid S)$ as a constant kernel $\delta_{w_0}$ independent of $S$. That
reading is degenerate: a constant kernel is Blackwell-equivalent to the null experiment, carries
zero information, and induces the trivial $\sigma$-algebra, which would make any garbling argument
vacuous. We avoid this by modelling each fixed source as a \emph{deterministic experiment that
observes a projection of the coordinate vector}~\eqref{eq:coordinates}.

\begin{definition}[Projection sources]
\label{def:projection}
Fix a partition of coordinate indices. A source $W$ that reveals the coordinates indexed by a
set $A \subseteq \{1,\dots,m\}$ is the deterministic kernel
\begin{equation}
\label{eq:projkernel}
\kappa_{W}(\cdot \mid S) \;=\; \delta_{\pi_A(S)}, \qquad \pi_A(S) := (S_j)_{j\in A},
\end{equation}
i.e.\ supplying $W$ lets a decision-maker condition on $\pi_A(S)$ and on nothing outside $A$.
Write $\Fcal_W := \sigma(\pi_A)$ for the sub-$\sigma$-algebra of $\Scal$ generated by
the projection. In our instantiation $W_1$ reveals coordinate $S_1$ (the call contract,
$A_1=\{1\}$), $W_2$ reveals $S_2$ (the wire invariant, $A_2=\{2\}$), and the superset
$W_{1+} = W_2 \,{+}{+}\, W_1$ reveals both ($A_{1+}=\{1,2\}$).
\end{definition}

Under Definition~\ref{def:projection} a fixed string is \emph{not} an uninformative constant: it
is a deterministic experiment whose informativeness is exactly the resolving power of its
$\sigma$-algebra $\Fcal_W$. The signal alphabet is the range of the projection, and the
``text varies with the task'' precisely because the realized coordinate values $\pi_A(S_T)$ differ
across task instances $T$. This is the well-posed reading that makes the comparison of
\S\ref{sec:order} a comparison of experiments rather than a metaphor; it is the construction the
internal due-diligence record flagged as required and that the earlier draft did not
state~\citep{torgersen1991}.

\subsection{The conditioning parameter \texorpdfstring{$\theta$}{theta}}
\label{subsec:theta}

The model rule $\delta_M$ carries an implicit parameter $\theta$: the model's weights and decoding
configuration, equivalently its \emph{prior over conventions} acquired in pretraining. We write
$I(S;W\mid\theta)$ for the model-relative value of a private source, the information $W$ supplies
\emph{beyond what $\theta$ already encodes}. When $\theta$ already ``knows'' a convention
(the model can guess it without the source), $W$'s marginal value collapses; when the convention is
private to $\theta$, $W$ is maximally valuable. This $\theta$-dependence is not a nuisance: it is
the formal home of the \emph{model-relative value of private context} measured in
\S\ref{sec:experiments}, where the guessability of a fixed convention spans $0\%$ to $75\%$ across
vendors. Throughout, all realized-PASS quantities are implicitly indexed by $\theta$; we suppress
it except where the model-relativity is the point.

% ==========================================================================
\section{The Blackwell order over context sources}
\label{sec:order}
% ==========================================================================

\subsection{The order and its two registers}

\begin{definition}[Garbling order]
\label{def:blackwell}
$W_1 \Bord W_2$ ($W_1$ \emph{Blackwell-dominates} $W_2$) iff there is a Markov kernel
(garbling) $G : \Wcal_1 \to \Delta(\Wcal_2)$ with $\kappa_{W_2} = G \circ \kappa_{W_1}$;
i.e.\ $W_2$'s signal is a noised, post-processed version of $W_1$'s, computable from $W_1$'s signal
alone without further access to $S$.
\end{definition}

The classical content of the order is the Blackwell--Sherman--Stein theorem, which we state with
hypotheses for a self-contained treatment.

\begin{theorem}[Blackwell--Sherman--Stein; universal dominance for the optimal rule]
\label{thm:blackwell}
Let $W_1, W_2$ be experiments on a common parameter space $\Scal$ with standard Borel signal
spaces. The following are equivalent:
\begin{enumerate}
  \item[\textup{(i)}] $W_1 \Bord W_2$ (Definition~\ref{def:blackwell});
  \item[\textup{(ii)}] for every prior $\mu$ on $\Scal$, every finite decision space, and every
    bounded loss $\ell$, the Bayes risk under $W_1$ is no larger than under $W_2$; equivalently the
    optimal value satisfies $V^\star(W_1,T) \ge V^\star(W_2,T)$ for every task $T$ with bounded
    utility.
\end{enumerate}
Randomized decision rules are admitted on both sides; if the dominated experiment is realized with
a randomized garbling, no loss of generality is incurred by allowing randomized rules in
\textup{(ii)}.
\end{theorem}

\begin{proof}
This is Blackwell's theorem~\citep{blackwell1953}; the loss-function form and the role of
randomization are standard~\citep{lecam1986,torgersen1991}. We use it as a black box. The only
substantive specialization is the reading of ``decision problem'' as ``task with a bounded
verifier utility'' and ``experiment'' as ``context source'' (Definition~\ref{def:source}), which
is faithful because $v_T \in \{0,1\}$ is a bounded utility and the kernels of
Definition~\ref{def:projection} are standard Borel.
\end{proof}

\begin{corollary}[Universal context dominance, C1]
\label{cor:c1}
If $W_1 \Bord W_2$ then, for the \emph{optimal} user, supplying $W_1$ is at least as good as
supplying $W_2$ on \emph{every} task, prior, and bounded utility, a decision-independent,
for-all-tasks guarantee that no task-specific scalar score (relevance, mutual information,
$\Vcal$-usable information) provides.
\end{corollary}

\begin{remark}[The order is partial; no lattice]
\label{rem:partial}
$\Bord$ is a genuine partial order: there exist pairs $W_1, W_2$ with neither $W_1 \Bord
W_2$ nor $W_2 \Bord W_1$, and the order does not form a lattice~\citep{torgersen1991}.
Incomparability is exactly the regime our experiments target.
\end{remark}

\subsection{Deterministic sources: incomparability via \texorpdfstring{$\sigma$}{sigma}-algebra non-refinement}

For deterministic projection sources (Definition~\ref{def:projection}) the garbling order reduces
to refinement of the generated $\sigma$-algebras, which makes incomparability checkable by a
purely structural argument rather than by an empirical divergence.

\begin{lemma}[Refinement characterization for projection sources]
\label{lem:noref}
Let $W_A, W_B$ be projection sources revealing coordinate sets $A,B$
(Definition~\ref{def:projection}), with coordinates $S_1,\dots,S_m$ mutually independent under
some full-support prior and each $S_j$ non-degenerate. Then:
\begin{enumerate}
  \item[\textup{(a)}] $W_A \Bord W_B \iff \Fcal_{W_B} \subseteq \Fcal_{W_A}
    \iff B \subseteq A$;
  \item[\textup{(b)}] consequently $W_A$ and $W_B$ are Blackwell-\emph{incomparable} iff
    $A \not\subseteq B$ and $B \not\subseteq A$, i.e.\ each reveals at least one coordinate the
    other hides.
\end{enumerate}
\end{lemma}

\begin{proof}
\emph{(a)} For deterministic experiments a garbling $G$ with $\kappa_{W_B}=G\circ\kappa_{W_A}$
must compute $\pi_B(S)$ from $\pi_A(S)$ alone (no further access to $S$). Such a measurable map
exists iff $\pi_B$ is $\sigma(\pi_A)$-measurable, i.e.\ $\Fcal_{W_B}\subseteq\Fcal_{W_A}$.
Under coordinate independence and non-degeneracy, $\sigma(\pi_B)\subseteq\sigma(\pi_A)$ holds iff
every coordinate in $B$ is among those in $A$, i.e.\ $B\subseteq A$: if some $j\in B\setminus A$,
then $S_j$ is independent of $\pi_A(S)$ and non-constant, so it is not
$\sigma(\pi_A)$-measurable, and no garbling can reconstruct it; conversely if $B\subseteq A$ the
coordinate-deletion map is a (deterministic) garbling. \emph{(b)} is the contrapositive of (a)
applied in both directions.
\end{proof}

Lemma~\ref{lem:noref} is the formal version of the garbling argument used in
\S\ref{sec:experiments}: from the wire format ($W_2$, revealing $S_2$) one cannot manufacture the
private arg-order / required-keyword / exception-type facts ($S_1$), and from the call contract
($W_1$, revealing $S_1$) one cannot manufacture the cents-scale / check-digit facts ($S_2$). With
$A_1=\{1\}$, $A_2=\{2\}$ neither set contains the other, so $W_1, W_2$ are Blackwell-incomparable
by construction (and $W_{1+}$, revealing $\{1,2\}$, dominates both). This guarantees the
\emph{theoretical} incomparability the experiments then exhibit empirically.

\subsection{Bridging the optimal-rule order to the fixed model}
\label{subsec:bridge}

Theorem~\ref{thm:blackwell} concerns $V^\star$. Our experiments measure $\PASS_M$ for a
single fixed, generally sub-optimal rule $\delta_M$. The two are not interchangeable: an inequality
between optimal values need not hold for a fixed rule, and, as we show empirically, the fixed rule
can even \emph{violate} monotonicity that the optimal value obeys. We therefore introduce a
$\delta_M$-restricted order and state precisely what it does and does not inherit from
Theorem~\ref{thm:blackwell}.

\begin{definition}[$\delta_M$-restricted, $\Dcal$-restricted order]
\label{def:dmorder}
Fix the model rule $\delta_M$ and a finite task set $\Dcal$. Say $W_1 \succeq_{M,\Dcal} W_2$ iff
$\PASS_M(W_1,T) \ge \PASS_M(W_2,T)$ for every $T \in \Dcal$. Call $W_1, W_2$
\emph{$(M,\Dcal)$-incomparable} iff $W_1 \not\succeq_{M,\Dcal} W_2$ and $W_2 \not\succeq_{M,\Dcal} W_1$, i.e.\
each strictly out-performs the other on some task in $\Dcal$.
\end{definition}

$\succeq_{M,\Dcal}$ is the order our certificate measures. It is decision-restricted (one fixed rule)
and distribution-restricted (a fixed finite $\Dcal$); within those restrictions it is still a genuine,
\emph{for-all-tasks-in-$\Dcal$} dominance statement, not a scalar.

\begin{lemma}[One-way bridge]
\label{lem:bridge}
$(M,\Dcal)$-incomparability is strictly weaker than Blackwell incomparability and does \emph{not}
follow from it without a value-attainment hypothesis:
\begin{enumerate}
  \item[\textup{(a)}] If $W_1, W_2$ are $(M,\Dcal)$-incomparable for some $\delta_M$ and $\Dcal$, then they
    are Blackwell-incomparable. \textup{(Empirical incomparability implies structural
    incomparability.)}
  \item[\textup{(b)}] The converse fails in general: Blackwell incomparability concerns $V^\star$,
    and a fixed sub-optimal $\delta_M$ may fail to realize the optimal-value crossover on $\Dcal$.
    It holds under the value-attainment hypothesis (Assumption~\ref{ass:attain}).
\end{enumerate}
\end{lemma}

\begin{assumption}[Value attainment on the witnessing tasks]
\label{ass:attain}
There exist tasks $T_a, T_b$ and a margin $\gamma>0$ such that $\delta_M$ attains
$\PASS_M(W_1,T_a) - \PASS_M(W_2,T_a) \ge \gamma$ and
$\PASS_M(W_2,T_b) - \PASS_M(W_1,T_b) \ge \gamma$; i.e.\ the fixed model actually
exploits each source's private coordinate on the task that requires it. Under
Assumption~\ref{ass:attain} with $\{T_a,T_b\}\subseteq \Dcal$, $W_1$ and $W_2$ are
$(M,\Dcal)$-incomparable.
\end{assumption}

\begin{proof}[Proof of Lemma~\ref{lem:bridge}]
\emph{(a)} $(M,\Dcal)$-incomparability exhibits a task $T_b$ where
$\PASS_M(W_2,T_b) > \PASS_M(W_1,T_b)$. Were $W_1 \Bord W_2$, then $W_2$'s signal
is a garbling of $W_1$'s, so the rule ``garble $W_1$'s signal to $W_2$'s, then apply $\delta_M$''
is available under $W_1$ and achieves $\PASS_M(W_2,T_b)$ using $W_1$; hence $W_1$ admits a
rule matching $W_2$ on $T_b$, while $\delta_M$ itself does strictly worse, no contradiction with
$W_1\Bord W_2$ by itself, but the symmetric task forbids the reverse garbling. Concretely, a task
$T_a$ with the reverse strict inequality $\PASS_M(W_1,T_a) > \PASS_M(W_2,T_a)$ forbids
$W_2 \Bord W_1$ by the same construction with roles exchanged. Since $(M,\Dcal)$-incomparability
supplies both $T_a$ and $T_b$, neither garbling direction can hold, so $W_1,W_2$ are
Blackwell-incomparable.\footnote{The argument uses only that a garbling is post-processing of the
signal, available to \emph{any} fixed rule; it does not require optimality of $\delta_M$. For the
deterministic projection sources of Definition~\ref{def:projection} the garbled rule is admissible
without measure-theoretic subtlety, since the garbling is a measurable coordinate map.} \emph{(b)}
Assumption~\ref{ass:attain} directly yields $\PASS_M(W_1,T_a) > \PASS_M(W_2,T_a)$ and the reverse
on $T_b$, which is Definition~\ref{def:dmorder}'s incomparability; the failure of the unconditional
converse is the content of \S\ref{subsec:antimono} (anti-monotonicity: a fixed rule violating an
optimal-value inequality).
\end{proof}

Lemma~\ref{lem:bridge} is the keystone fix. It says the empirical program is sound in the direction
we actually need: \emph{a certified $(M,\Dcal)$-incomparability (which our certificate delivers)
implies genuine Blackwell incomparability} (part (a)), so the structural claim is earned, not
assumed; and it makes explicit that the reverse implication, predicting the realized crossover
from the abstract theorem, requires Assumption~\ref{ass:attain}, which our task design enforces
(each cell is constructed so the decisive private coordinate is necessary for PASS) and our
measurements verify.

% ==========================================================================
\section{The incomparability theorem}
\label{sec:incomparability}
% ==========================================================================

We now establish that no scalar can rank Blackwell-incomparable sources. We \emph{split} the
claim into a clean theorem for query-\emph{independent} source scores $\varphi(W)$ and a separate,
explicitly empirical statement for query-\emph{conditioned} scores $\varphi(W,T)$ (the relevance
score practitioners actually use); an earlier draft conflated the two and over-claimed.

\subsection{Query-independent scores: the theorem}

\begin{theorem}[No query-independent scalar ranks incomparable sources, C2]
\label{thm:incomparable}
Let $W_1, W_2$ be Blackwell-incomparable. Let $\varphi : \{\text{sources}\} \to \mathbb R$ be any
\emph{query-independent} scalar score (a single number per source: relevance to a fixed corpus,
mutual information $I(S;W)$, $\Vcal$-usable information, perplexity-gap, expected information gain
under a fixed prior). Then there exist tasks $T_a, T_b$ such that the ranking of $\{W_1,W_2\}$
induced by $\varphi$ disagrees with the optimal-value ranking on at least one of $T_a, T_b$; and,
under Assumption~\ref{ass:attain}, with the realized-$\PASS_M$ ranking on at least one of them.
\end{theorem}

\begin{proof}
By incomparability, neither source garbles the other. By the contrapositive of
Theorem~\ref{thm:blackwell}, ``not $W_1 \Bord W_2$'' yields a task $T_b$ (a prior plus bounded
utility) with $V^\star(W_2,T_b) > V^\star(W_1,T_b)$, and ``not $W_2 \Bord W_1$'' yields a task
$T_a$ with $V^\star(W_1,T_a) > V^\star(W_2,T_a)$; both inequalities are \emph{strict}, so ties do
not arise. A query-independent $\varphi$ assigns the single ordered pair $(\varphi(W_1),
\varphi(W_2))$, fixing one total order of $\{W_1,W_2\}$ \emph{independent of the task}. That fixed
order agrees with the optimal-value order on at most one of $T_a, T_b$ (they point opposite ways),
hence is violated on the other. The $\PASS_M$ conclusion follows by replacing $V^\star$ with
$\PASS_M$ via Assumption~\ref{ass:attain}, which guarantees the realized rule reproduces both
strict inequalities with margin $\gamma$.
\end{proof}

\begin{remark}[Scope]
\label{rem:c2scope}
Theorem~\ref{thm:incomparable} rules out exactly the class of \emph{source-level} scalars: any
$\varphi$ that is a fixed number per source, regardless of task. $\Vcal$-usable information, mutual
information, perplexity-gap, and (fixed-prior) EIG all live in this class, so they are corollary
victims of the theorem, not competitors (Proposition~\ref{prop:vinfo-monotone}). The theorem is
mathematically a one-line consequence of Blackwell's theorem; the contribution is the
\emph{application and framing}, reading cross-task non-transfer of relevance as the signature of a
partial order, not new core mathematics~\citep{blackwell1953,torgersen1991}.
\end{remark}

\begin{proposition}[$\Vcal$-usable information is Blackwell-monotone]
\label{prop:vinfo-monotone}
Let $\Vcal$-usable information $I_{\Vcal}(W\!\to\! S)$ be defined, after
\citet{ethayarajh2022}, as the reduction in expected predictive Bayes risk of the best predictor
in a fixed family $\Vcal$ when given the signal $w$ versus the null signal. If $W_1 \Bord W_2$
and the predictor family $\Vcal$ is closed under post-composition with measurable garblings, then
$I_{\Vcal}(W_1\!\to\! S) \ge I_{\Vcal}(W_2\!\to\! S)$. Hence $I_{\Vcal}$ is a Blackwell-monotone
scalar: it respects $\Bord$ but collapses it to one real number.
\end{proposition}

\begin{proof}
Let $G$ be the garbling with $\kappa_{W_2}=G\circ\kappa_{W_1}$. Any predictor $g\in\Vcal$ acting
on $W_2$'s signal is realized on $W_1$'s signal by the composite $g\circ G$, which lies in $\Vcal$
by the closure hypothesis. Therefore the best $\Vcal$-predictor under $W_1$ attains at least the
predictive accuracy of the best under $W_2$, so its expected predictive risk is no larger; the
risk-reduction relative to the common null signal is thus at least as large under $W_1$. The
scalarity is immediate: $I_{\Vcal}$ returns a single real number per source, so it induces a total
order and, by Theorem~\ref{thm:incomparable}, cannot rank a Blackwell-incomparable pair correctly
for all tasks.
\end{proof}

\begin{remark}[Scope of Proposition~\ref{prop:vinfo-monotone}]
\label{rem:vinfo-scope}
The closure hypothesis (a family closed under post-composition with garblings) is the natural
sufficient condition and holds for predictor families rich enough to apply a fixed measurable
post-processing; for narrower families monotonicity may fail and we do not claim it. We prove
monotonicity for the closest neighbour, $\Vcal$-usable information, and conjecture, rather than
prove, the analogue for mutual information, perplexity-gap, and fixed-prior EIG, each of which is
a data-processing-monotone functional of the same kernels.
\end{remark}

\subsection{Query-conditioned scores: the relevance trap (empirical)}
\label{subsec:trap}

A query-conditioned score $\varphi(W,T)$, e.g.\ the lexical or embedding relevance of a source to
the task prompt, which is what retrieval and RAG actually deploy, induces a \emph{different} order
per task and is therefore \emph{not} a single fixed total order. Theorem~\ref{thm:incomparable}
does \emph{not} cover it. We make two precise statements instead.

\begin{proposition}[Constrained impossibility for feature-measurable scores]
\label{prop:trap-constrained}
Suppose $\varphi(W,T)$ is measurable with respect to a fixed feature map of the pair $(W,T)$ that
is \emph{topical}: it depends on $W,T$ only through surface lexical overlap and is invariant to the
private coordinate $\pi_A(S_T)$ that the verifier checks. If there is a task $T_c$ whose decisive
coordinate is supplied by $W_1$ but whose surface lexicon is saturated with $W_2$'s vocabulary,
then $\varphi(W_2,T_c) > \varphi(W_1,T_c)$ while the realized utility ranks $W_1 \succ W_2$ on
$T_c$. Hence no topical query-conditioned score ranks $\{W_1,W_2\}$ correctly on all tasks.
\end{proposition}

\begin{proof}
By topical invariance, $\varphi(\cdot,T_c)$ ignores the verifier-decisive coordinate and orders by
lexical overlap, which favours $W_2$ on $T_c$ by construction; utility favours $W_1$ because $T_c$'s
PASS requires $W_1$'s coordinate. The two orders disagree on $T_c$.
\end{proof}

Proposition~\ref{prop:trap-constrained} is the strongest \emph{general} statement we can make about
query-conditioned scores: it holds for any score in the topical (lexical-feature) class, but not for
an arbitrary learned reranker that could, in principle, read the private coordinate. We therefore
treat the defeat of query-conditioned relevance as an \emph{empirical} result, demonstrated for the
concrete lexical-cosine score practitioners use. In \S\ref{sec:experiments} the
\texttt{trap\_store\_wire} task realizes $T_c$: bag-of-words cosine gives $\varphi(W_1)=0.36 <
\varphi(W_2)=0.44$ (ranking $W_2\succ W_1$), yet realized PASS ranks $W_1$ at $96\%$--$100\%$ versus
$W_2$ at $0\%$ across all six models. The claim is restricted to lexical relevance; we do not claim
to have defeated every learned reranker.

% ==========================================================================
\section{The decision-restricted deficiency surrogate and its uniform certificate}
\label{sec:estimator}
% ==========================================================================

\subsection{The surrogate and its correct name}
\label{subsec:estimator-name}

Exact garblings over text are uncomputable and $\Bord$ is partial, so we measure the
$\delta_M$-restricted order of Definition~\ref{def:dmorder}. Fix the model rule $\delta_M$ and a
finite task set $\Dcal$ and define the \emph{decision-restricted value-deficiency} (a one-sided
worst-case regret of $W_1$ relative to $W_2$ on $\Dcal$):
\begin{equation}
\label{eq:estimator}
\dhat_{\Dcal}(W_1,W_2) \;:=\; \max\!\Big(0,\ \sup_{T\in \Dcal}\big[\PASS_M(W_2,T)-\PASS_M(W_1,T)\big]\Big),
\end{equation}
estimated per cell by $\PASS_M(W,T) \approx k_{W,T}/n$ from $n$ completions. Then
$\dhat_{\Dcal}(W_1,W_2)=0$ certifies empirical dominance $W_1 \succeq_{M,\Dcal} W_2$, and mutual
positivity $\dhat_{\Dcal}(W_1,W_2)>0 \wedge \dhat_{\Dcal}(W_2,W_1)>0$ certifies
$(M,\Dcal)$-incomparability.

\begin{remark}[This is \emph{not} Le~Cam deficiency]
\label{rem:notlecam}
The Le~Cam deficiency $\delta(\mathcal E,\mathcal F) = \inf_{G}\ \sup_{\text{loss}} \big[
R_{\mathcal F} - R_{G\circ\mathcal E}\big]$ \emph{minimizes over garblings} $G$ and takes a
\emph{supremum over a class of losses}~\citep{lecam1964,lecam1986,torgersen1991}. The quantity in
\eqref{eq:estimator} performs no garbling minimization and uses a single fixed rule $\delta_M$ over
a fixed finite task set; it is a one-sided value/regret gap, not a deficiency. No general inequality
relates $\dhat_{\Dcal}$ to the true Le~Cam deficiency for a fixed sub-optimal rule (it can be
larger or smaller in either direction). We therefore name \eqref{eq:estimator} the
\emph{decision-restricted value-deficiency} and use ``Le~Cam deficiency'' only as the conceptual
analogy that motivates a graded relaxation of the partial order. The one direction that \emph{does}
transfer is qualitative: $\dhat_{\Dcal}(W_1,W_2)=0$ realizes the for-all-tasks-in-$\Dcal$ dominance
of Definition~\ref{def:dmorder}, the measurable shadow of Corollary~\ref{cor:c1}.
\end{remark}

\subsection{The uniform certificate}
\label{subsec:certificate}

The surrogate \eqref{eq:estimator} selects the \emph{argmax-gap} task, so a pointwise confidence
statement on a single task ignores the selection. We pay the multiplicity with an exact
Clopper--Pearson interval per cell and a union bound, which is distribution-free and finite-sample
valid.

\begin{proposition}[Uniform incomparability certificate]
\label{prop:certificate}
Within each cell $(W,T)$ draw $n$ i.i.d.\ Bernoulli PASS outcomes (see
Assumption~\ref{ass:iid}). Let $[L_{W,T},U_{W,T}]$ be the exact (two-sided) Clopper--Pearson
interval~\citep{clopperpearson1934} for $p_{W,T}=\PASS_M(W,T)$ at level
$\alpha = \eta/(2|\Dcal|)$. By the Bonferroni/union bound~\citep{bonferroni1936} over the $2|\Dcal|$
proportions $\{p_{W_1,T},p_{W_2,T}\}_{T\in \Dcal}$, all $2|\Dcal|$ intervals hold simultaneously with
probability at least $1-\eta$. On that event,
\begin{equation}
\label{eq:LD}
L_{\Dcal}(W_1,W_2) \;:=\; \max_{T\in \Dcal}\big[\,L_{W_2,T} - U_{W_1,T}\,\big]
\end{equation}
is a $(1-\eta)$ lower confidence bound on $\sup_{T\in \Dcal}[p_{W_2,T}-p_{W_1,T}]$. Hence
$L_{\Dcal}(W_1,W_2) > 0$ \textbf{certifies} $\dhat_{\Dcal}(W_1,W_2)>0$ at confidence $\ge 1-\eta$.
Both directions $L_{\Dcal}(W_1,W_2)$ and $L_{\Dcal}(W_2,W_1)$ read off the \emph{same} $2|\Dcal|$ intervals, so the
joint incomparability claim $\{\dhat_{\Dcal}(W_1,W_2)>0 \wedge \dhat_{\Dcal}(W_2,W_1)>0\}$
holds at $\ge 1-\eta$ with no additional correction.
\end{proposition}

\begin{proof}
Let $E$ be the event that all $2|\Dcal|$ Clopper--Pearson intervals contain their true proportions.
Each interval fails with probability at most $\alpha=\eta/(2|\Dcal|)$ (exact coverage of
Clopper--Pearson~\citep{clopperpearson1934}), so by the union bound $\Pr[E^c]\le 2|\Dcal|\cdot\alpha =
\eta$, giving $\Pr[E]\ge 1-\eta$. On $E$, for every $T$ we have $p_{W_2,T}\ge L_{W_2,T}$ and
$p_{W_1,T}\le U_{W_1,T}$, so $p_{W_2,T}-p_{W_1,T} \ge L_{W_2,T}-U_{W_1,T}$; taking $\max_T$ on both
sides, $\sup_T[p_{W_2,T}-p_{W_1,T}] \ge L_{\Dcal}(W_1,W_2)$. Thus on $E$ (probability $\ge 1-\eta$) the
true sup is at least $L_{\Dcal}(W_1,W_2)$, which is the claimed lower confidence bound; if it is positive,
the true sup is positive, certifying $\dhat_{\Dcal}(W_1,W_2)>0$. The reverse direction uses
$L_{W_1,T}$ and $U_{W_2,T}$, which are endpoints of the \emph{same} $2|\Dcal|$ intervals already
included in $E$; therefore both one-sided certificates are valid on the single event $E$, and their
conjunction holds at $\ge 1-\eta$ without a further Bonferroni factor.
\end{proof}

\begin{remark}[Why a naive DKW bound fails]
\label{rem:dkw}
The Dvoretzky--Kiefer--Wolfowitz inequality~\citep{dvoretzky1956} bounds the sup-deviation of
\emph{one} empirical CDF from \emph{its own} population CDF over a single i.i.d.\ sample. Our object
is different in two ways that DKW does not handle: it is a $\max$ over $|\Dcal|$ \emph{differences of
two binomial means across distinct cells}, and the $\max$ \emph{selects} the argmax-gap cell. A
single-sample sup inequality neither spans distinct cells nor pays for the selection, so its nominal
level under-covers a supremum over $|\Dcal|$ comparisons; the Type-I error grows with $|\Dcal|$. The
multiplicity must be paid explicitly via the union bound (Proposition~\ref{prop:certificate}), not
absorbed into a single-sample DKW envelope. This is the correction over the earlier sketch.
\end{remark}

\begin{assumption}[Sampling regime]
\label{ass:iid}
The $n$ outcomes within a cell are i.i.d.\ Bernoulli draws of $v_T(\delta_M(T,w))$. This holds when
completions are sampled independently at fixed temperature/top-$p$ with independent randomness per
draw; for the single-shot arms (one completion + deterministic code extraction per draw) the draws
are manifestly independent. For the agentic arms (\texttt{claude -p}, multi-turn) independence is
across \emph{trajectories}: each of the $n$ draws is a fresh trajectory with no shared state, so the
binomial model applies to the trajectory-level PASS outcome; we report decoding configuration and
defend this independence in \S\ref{sec:experiments}. Where draws are near-deterministic
($p\in\{0,1\}$ stark cells), Clopper--Pearson collapses and the certificate is conservative.
\end{assumption}

\subsection{The interference certificate (anti-monotonicity)}
\label{subsec:interference}

The deficiency surrogate also detects a phenomenon the optimal-value order forbids but a fixed rule
can exhibit: supplying a \emph{superset} source $W_{1+}$ (revealing strictly more coordinates than
$W_1$; Definition~\ref{def:projection}) can \emph{degrade} performance. Define the per-task
interference
\begin{equation}
\label{eq:interference}
\Psi_T \;:=\; \PASS_M(W_{1+},T) - \PASS_M(W_1,T) - \PASS_M(W_2,T) + \PASS_M(\mathrm{none},T),
\end{equation}
the second-difference (interaction) of adding $W_2$ on top of $W_1$. We write the interaction as
$\Psi_T$ (rather than a symbol $I$ that would collide with mutual information $I(S;W\mid\theta)$). A
certified $\Psi_T<0$ on a cell whose $W_1$ arm is at ceiling is the localized statement that ``more
strictly-more-informative context hurts.''

\begin{proposition}[Interference certificate]
\label{prop:interference}
Fix a target task $T^\star\in \Dcal$ chosen \emph{before} unblinding. With per-cell exact
Clopper--Pearson endpoints $\mathrm{CP}_{\mathrm{lo}},\mathrm{CP}_{\mathrm{hi}}$ at level
$\alpha=\eta/4$ on the four arms $\{W_{1+},W_1,W_2,\mathrm{none}\}$ of $T^\star$,
\begin{equation}
\label{eq:psiupper}
\overline\Psi_{T^\star} \;:=\; \mathrm{CP}_{\mathrm{hi}}(p_{W_{1+}}) - \mathrm{CP}_{\mathrm{lo}}(p_{W_1}) - \mathrm{CP}_{\mathrm{lo}}(p_{W_2}) + \mathrm{CP}_{\mathrm{hi}}(p_{\mathrm{none}})
\end{equation}
is a $(1-\eta)$ upper confidence bound on $\Psi_{T^\star}$. Hence $\overline\Psi_{T^\star}\le-\tau$
certifies harmful non-monotonicity $\Psi_{T^\star}\le-\tau$ at confidence $\ge 1-\eta$. If instead
the harmful cell is \emph{selected} from $|\Dcal|$ candidates after unblinding, use $\alpha=\eta/(4|\Dcal|)$;
this is the same selection correction the certificate applies to the deficiency surrogate
(Remark~\ref{rem:dkw}), and omitting it would repeat the multiplicity error we criticize.
\end{proposition}

\begin{proof}
The four arms are independent Bernoulli samples (separate completions), so on the intersection of
the four Clopper--Pearson events, probability $\ge 1-4\alpha = 1-\eta$ by the union
bound, each true proportion lies in its interval. Substituting the coverage-respecting endpoint for
each term of \eqref{eq:interference} (upper endpoints for the positively-signed $p_{W_{1+}},
p_{\mathrm{none}}$, lower endpoints for the negatively-signed $p_{W_1}, p_{W_2}$) gives
$\Psi_{T^\star}\le\overline\Psi_{T^\star}$ on that event. For a post-hoc selected cell, replacing
$\alpha$ by $\eta/(4|\Dcal|)$ and union-bounding over the $4|\Dcal|$ endpoints preserves $\ge 1-\eta$
simultaneous coverage.
\end{proof}

\begin{remark}[The leaky-baseline limitation, stated as a protocol]
\label{rem:leaky}
The upper bound \eqref{eq:psiupper} includes $+\mathrm{CP}_{\mathrm{hi}}(p_{\mathrm{none}})$, so it
can certify harm only on cells whose \texttt{none}-baseline floors near $0$ (the convention is
unguessable for the model under test). When the baseline is non-zero, e.g.\ \texttt{api\_argorder}
at $\texttt{none}=75\%$ for GPT-5.5, $45\%$ for Kimi, $15\%$ for Opus, the bound is mechanically
slack and harm is not certifiable even when the point estimate is strongly negative. We therefore
adopt the explicit protocol: \emph{certify $\Psi_T$ only on cells whose \texttt{none}-baseline
floors near $0$ for the model under test}, and report, per model, which cell carried the certificate
and what its baseline was. This makes the cell selection a stated rule rather than an outcome-driven
choice. The same model-relative guessability that defines the value of private context
(\S\ref{subsec:theta}) thus governs which cells can carry the interference certificate.
\end{remark}

\subsection{Anti-monotonicity as information non-monotonicity of \texorpdfstring{$\delta_M$}{the model rule}}
\label{subsec:antimono}

A certified $\Psi_{T^\star}<0$ on a ceiling-$W_1$ cell is more than a curiosity: it is a falsification
of Blackwell superset-dominance \emph{for the fixed model rule}, and it pinpoints exactly where the
optimal-rule order (Theorem~\ref{thm:blackwell}) and the realized rule diverge.

\begin{theorem}[Information non-monotonicity of the model rule]
\label{thm:antimono}
Let $W_{1+}$ reveal a superset of $W_1$'s coordinates, so $W_{1+}\Bord W_1$
(Lemma~\ref{lem:noref}). Any \emph{Blackwell-coherent} rule, one whose value is monotone under
$\Bord$, as the Bayes rule $\delta^\star$ is by Theorem~\ref{thm:blackwell}, satisfies
$\mathrm{value}(W_{1+},T)\ge\mathrm{value}(W_1,T)$ for every task $T$. If a fixed model rule
$\delta_M$ has a task $T^\star$ with $\PASS_M(W_{1+},T^\star) <
\PASS_M(W_1,T^\star)$, then $\delta_M$ is \emph{not} Blackwell-coherent: it violates
monotonicity under the garbling order on $T^\star$. In particular no monotone-submodular scalar
context-value score can predict $\delta_M$'s behaviour on $T^\star$.
\end{theorem}

\begin{proof}
$W_{1+}\Bord W_1$ by deletion of the extra coordinate (a deterministic garbling,
Lemma~\ref{lem:noref}(a)). By Theorem~\ref{thm:blackwell} any rule whose value respects $\Bord$
weakly improves under the dominating experiment. The hypothesized strict decrease
$\PASS_M(W_{1+},T^\star)<\PASS_M(W_1,T^\star)$ contradicts monotonicity, so
$\delta_M$ is not Blackwell-coherent. A monotone-submodular scalar $\phi$ assigns
$\phi(W_{1+})\ge\phi(W_1)$ (monotonicity), hence orders $W_{1+}$ above $W_1$ and cannot reproduce
the realized reversal.
\end{proof}

Theorem~\ref{thm:antimono} converts the empirical collapse into a structural statement: the harm is
not noise but a witness that real LLM decision rules are not Blackwell-coherent, which is precisely
the gap Lemma~\ref{lem:bridge}(b) warned cannot be closed without value attainment. The deficiency
surrogate detects this reversal because it tracks realized PASS, whereas every monotone scalar must
miss it by construction.

\subsection{Family-wise error across the certified claims}
\label{subsec:fwer}

Each certificate above controls error \emph{marginally} at $\eta$ for its own claim
(incomparability via Proposition~\ref{prop:certificate}; interference via
Proposition~\ref{prop:interference}). When several claims are read off the \emph{same} $n$-sample
run at the same $\eta$, incomparability (both directions, on one shared event $E$), one-directional
dominance, and the interference cells, the conjunction over distinct claims is \emph{not} jointly
controlled at $\eta$ unless the budget is split across claims. We state results as marginally
certified per claim and, where a joint statement is wanted, allocate $\eta$ across the $C$ certified
claims (e.g.\ $\eta/C$ each) so the global family-wise error is $\le\eta$. The incomparability pair
is the one exception that is genuinely joint at $\eta$ on a single event, as proved in
Proposition~\ref{prop:certificate}; all other cross-claim conjunctions on a shared run are reported
as marginal.

\subsection{Power and the role of \texorpdfstring{$n$}{n}}
\label{subsec:power}

The asymmetric sample sizes in \S\ref{sec:experiments} (Opus at $n=20$ certifies incomparability but
not interference) are principled, not ad hoc. For a stark cell with empirical proportion at an
endpoint ($\hat p\in\{0,1\}$), the Clopper--Pearson interval collapses to a one-sided bound that
clears $0$ at small $n$: the incomparability certificate \eqref{eq:LD} is carried by such stark
cells ($W_1$ at $\approx100\%$, $W_2$ at $0\%$ and vice versa), so $n=20$ suffices. The interference
certificate \eqref{eq:psiupper} mixes a possibly-intermediate \texttt{none} baseline and the
$W_{1+}$ arm; with non-endpoint proportions the intervals are wide and $n=20$ is underpowered. The
two certificates therefore have genuinely different effective-$n$ requirements, which is why the
consolidated record marks Opus incomparability as certified and Opus interference as reproduced but
not certified.

% ==========================================================================
\section{Deficiency-aware context selection}
\label{sec:method}
% ==========================================================================

The partial order is not only a diagnosis; it prescribes a selection rule. If sources form a
genuine partial order, the correct primitive is not ``score and take the top $k$'' but ``keep the
certified non-dominated set, then diversify over what remains.'' We give that rule, prove three
properties, exhibit it on our measured data without running a new experiment, and state its cost
and scope. The point is to convert the impossibility result of
\S\ref{sec:incomparability} and the anti-monotonicity of \S\ref{sec:estimator} into a constructive,
certificate-backed procedure.

\subsection{The certified empirical dominance relation}
Fix the model rule $\delta_M$ and a probe battery $\Dcal$. Write $W_i \trianglerighteq W_j$
(``$W_i$ empirically dominates $W_j$'') when the certificates of \S\ref{sec:estimator} jointly
establish, at confidence $\ge 1-\eta$, that (a) $\dhat_{\Dcal}(W_i,W_j)=0$ up to the equivalence
margin $\alpha_{\mathrm{tol}}$ (no probe task on which $W_j$ beats $W_i$ by more than the margin),
and (b) $L_{\Dcal}(W_j,W_i)>0$ ($W_i$ strictly beats $W_j$ on some probe task). This is exactly the
one-directional dominance the dominance control certifies. Sources that are mutually non-dominated
are \emph{certified incomparable} (both $L_{\Dcal}>0$); they form the empirical Pareto front.

\subsection{The selector}
\noindent\textbf{Algorithm 1 (\textsc{Probe-Select}).} \emph{Input:} candidate sources $\mathcal{C}$,
probe battery $\Dcal$, budget $k$, model $\delta_M$, level $\eta$.
\begin{enumerate}[leftmargin=2.2em,itemsep=2pt]
  \item Measure $\PASS(W,T)$ for all $W\in\mathcal{C}$, $T\in\Dcal$ ($n$ draws per cell); form the
    $2|\mathcal{C}|\!\cdot\!|\Dcal|$ Clopper--Pearson intervals once.
  \item Build $\trianglerighteq$ on $\mathcal{C}$ from those intervals
    (Propositions of \S\ref{sec:estimator}); no extra model calls.
  \item \emph{Prune} every $W$ that is certified-dominated by some $W'\in\mathcal{C}$.
  \item Let $P$ be the surviving certified non-dominated set (the empirical Pareto front).
  \item If $|P|\le k$ return $P$; else greedily add to the output the survivor with the largest
    certified marginal PASS gain on probe tasks not yet covered by the chosen set, breaking ties by
    smallest $\dhat_{\Dcal}$ to the chosen set, until $k$ are selected.
\end{enumerate}

\subsection{Guarantees}
\begin{proposition}[Monotone safety]\label{prop:safe}
\textsc{Probe-Select} never returns a source certified-dominated by an available source. Hence,
with probability $\ge 1-\eta$, it never discards a source that strictly wins a probe task in favor
of one that is nowhere better than it by more than $\alpha_{\mathrm{tol}}$.
\end{proposition}
\begin{proof}
Step 3 removes only certified-dominated sources, and step 5 selects from the survivors $P$. By the
dominance certificate (\S\ref{sec:estimator}), a pruned $W$ has, on the joint $(1-\eta)$ event, no
probe task on which it beats its dominator by more than $\alpha_{\mathrm{tol}}$; the dominator is
retained or itself Pareto-equivalent, so no strict probe winner is sacrificed for a uniformly
no-better source. \end{proof}

\begin{proposition}[Strict generalization of top-$k$]\label{prop:generalize}
If $\trianglerighteq$ totally orders $\mathcal{C}$ (every pair certified-comparable), the Pareto
front is a chain and \textsc{Probe-Select} returns its top $k$ elements, recovering scalar top-$k$
selection. The procedure therefore extends, rather than competes with, the scalar pipeline: it
differs only on the incomparable pairs the scalar provably mis-handles
(Theorem~\ref{thm:incomparable}).
\end{proposition}

\begin{proposition}[Coverage over incomparable coordinates]\label{prop:coverage}
On an incomparable survivor set the greedy step (5) maximizes, up to the standard
$(1-1/e)$ factor of monotone submodular coverage, the number of distinct probe tasks solved by the
selected set, i.e.\ the number of distinct latent coordinates covered, rather than retaining
sources that solve the same tasks. This is the partial-order-grounded analogue of
maximal-marginal-relevance diversity, with ``diversity'' defined by certified task coverage rather
than embedding distance.
\end{proposition}

\subsection{Worked illustration on the measured data}
\label{sec:method-illustration}
Consider $\mathcal{C}=\{W_1, W_2, W_{1+}\}$ with $W_{1+}=W_2 +\!\!+\, W_1$. The \emph{idealized}
Blackwell order would collapse this set to $\{W_{1+}\}$ alone: $W_{1+}$ literally contains both
signals, so $W_{1+}\Bord W_2$ \emph{and} $W_{1+}\Bord W_1$ ($W_1$ is the garbling that drops the
$W_2$ text), and a ``more signal is strictly better'' heuristic keeps only the superset. On a real
fixed model this is exactly wrong: the certified \emph{empirical} relation is different because of
the anti-monotonicity of \S\ref{sec:estimator}. From the measured pass rates (Table~\ref{tab:percell}),
$W_1$ beats $W_{1+}$ on the API cells (e.g.\ Haiku $98\%$ vs.\ $24\%$), so $W_{1+}\not\trianglerighteq
W_1$; $W_{1+}$ beats $W_1$ on the encoding cell; and both dominate $W_2$. \textsc{Probe-Select}
therefore prunes only $W_2$ and returns the front $\{W_1, W_{1+}\}$, whereas the superset heuristic
returns $\{W_{1+}\}$ and collapses on every API task. The selector converts the anti-monotonicity
finding into a selection-time safeguard: by ordering on the decision-restricted surrogate
$\dhat_{\Dcal}$ rather than on idealized signal content, it keeps the component source whose value
the superset destroys. No new measurement is required for this conclusion; it reads directly off the
certified $n{=}80$ Haiku cells.

\subsection{Cost and scope}
\label{sec:method-scope}
Building $\trianglerighteq$ on $m=|\mathcal{C}|$ candidates costs $m\!\cdot\!|\Dcal|\!\cdot\!n$
verifier-gated model calls (the pairwise certificates reuse a single per-(source,task) PASS table,
so the cost is linear in $m$, not quadratic), plus $O(m^2)$ interval comparisons. This is
\emph{curation scale}: selecting among a repository's handful of private-convention documents, or
curating a few-shot pool of tens to hundreds of exemplars. It is \emph{not} corpus-scale retrieval
over $10^6$ passages, and we do not claim it is. Two routes to scale, both future work, are
immediate: \textbf{(i)} amortize the probe by learning a deficiency predictor
$\widehat{\delta}_\varphi(W_i,W_j)$ from cheap features, converting the online
$m\!\cdot\!|\Dcal|\!\cdot\!n$ cost into offline training plus $O(m^2)$ inference; and \textbf{(ii)}
use a cheap scalar retriever to pre-filter a shortlist and run \textsc{Probe-Select} only as a
partial-order re-ranking head, where the incomparable-set size is small. We validate the certified
estimator that powers steps 1--3 in \S\ref{sec:experiments}; a corpus-scale evaluation of the full
selector against strong reranker baselines is the central open problem we identify in
\S\ref{sec:limitations}, not a result we report.

% ==========================================================================
\section{Experiments}
\label{sec:experiments}
% ==========================================================================

We instantiate the theory of \S\S\ref{sec:order}--\ref{sec:estimator} on real language models. The
design has three commitments. First, every reported number is an empirical PASS rate from live model
calls verified by an executable test harness; no value is modeled, simulated, or extrapolated.
Second, the two context sources are constructed to be \emph{Blackwell-incomparable by construction}
(\S\ref{sec:design}), so the incomparability theorem (Theorem~\ref{thm:incomparable}) has a
non-vacuous instance to bite on. Third, every certified claim is reported with the exact sample
size, the Clopper--Pearson endpoints, and the per-cell PASS counts from which it is derived
(\S\ref{sec:results}), so the certificate is auditable rather than asserted.

\subsection{Design: two incomparable context sources}
\label{sec:design}

We define two \emph{private, unguessable} conventions on the \emph{same} module \texttt{ledger},
matched on every nuisance axis (length, roughly three sentences, exactly one private convention
each, same module, same task shape). Because both are on-topic for any ``ledger'' task, any lexical,
embedding, mutual-information, or $\Vcal$-usable score treats them as equally relevant
candidates, the precondition for Theorem~\ref{thm:incomparable} to apply.

\begin{description}
  \item[$W_1$, API call contract.] \texttt{ledger.post(account, amount, *, memo)}: \texttt{account}
  first, \texttt{amount} second positional, \texttt{memo} a \emph{required keyword-only} argument;
  returns an integer entry id $\ge 1$; raises \texttt{ledger.PostError} (not
  \texttt{ValueError}/\texttt{KeyError}) on failure.
  \item[$W_2$, wire-encoding invariant.] Amounts serialize to \texttt{<sign><minor-unit-digits>\#<check>},
  where $\mathrm{check} = (\text{sum of digits}) \bmod 10$; e.g.\ \$12.30 $\mapsto$ \texttt{"+1230\#6"}
  and $-\$0.07 \mapsto$ \texttt{"-7\#7"}. Decoding rejects a wrong check digit.
  \item[$W_{1+}$, literal superset.] $W_{1+} = W_2 \mathbin{+\!\!+} W_1$, the concatenation of
  $W_2$'s signal with $W_1$'s. As a string superset of $W_2$, it Blackwell-dominates $W_2$
  ($W_{1+} \Bord W_2$); it is the anti-monotonicity / dominance probe. We also use
  $W_{1+}^{\mathrm{rev}} = W_1 \mathbin{+\!\!+} W_2$ ($W_1$ first) as an order control.
\end{description}

\paragraph{Why incomparable (garbling argument).}
A garbling acts on the \emph{signal alone}. From $W_2$'s wire format one cannot manufacture the
private argument order, the required keyword-only \texttt{memo}, or the exception type; from $W_1$'s
contract one cannot manufacture the cents scale, the \texttt{\#} separator, or the mod-10 check. The
two facts are logically independent coordinates of the latent requirement $S$, so neither signal's
$\sigma$-algebra refines the other's, neither sufficiency direction holds, and $W_1, W_2$ are
Blackwell-incomparable. This is the concrete instance of the coordinate/projection construction of
\S\ref{sec:formalization}: $W_1$ and $W_2$ are deterministic experiments observing distinct,
mutually non-refining projections of $S$ (Lemma~\ref{lem:noref}).

\paragraph{Confound defense: the double crossover.}
The experiment's signature is a \emph{double crossover} measured against a \texttt{none} baseline (no
convention supplied): $W_1$ lifts the API tasks but \emph{not} the encoding task, while $W_2$ lifts
the encoding task but \emph{not} the API tasks. If the encoding family were merely ``harder,'' $W_1$
would not lift it and $W_2$ would still have to lift its own family from the same floor. Matched
\texttt{none} baselines that floor near zero on both families therefore establish equal intrinsic
hardness, so any divergence is attributable to the matching private fact, not to a difficulty
asymmetry.

\subsection{Tasks}
\label{sec:tasks}

The full design comprises seven tasks with stub-\texttt{ledger} verifiers (return code $0 =$ PASS):
four $W_1$-decisive API tasks (\texttt{api\_post\_ok}, \texttt{api\_argorder}, \texttt{api\_error\_type},
\texttt{api\_memo\_required}), two $W_2$-decisive encoding tasks (\texttt{enc\_amount}, \texttt{dec\_amount}),
and one relevance trap (\texttt{trap\_store\_wire}). The verifier stubs a fake \texttt{ledger} that
accepts \emph{only} the exact private contract (rejecting wrong positional order, a missing
\texttt{memo=}, or the wrong exception type).

\paragraph{Certified battery.}
The certified runs use a fixed four-cell battery: \texttt{api\_post\_ok}, \texttt{api\_argorder}
(hardened so \texttt{memo} arrives through a parameter, forcing reliance on the keyword-only fact),
\texttt{enc\_amount}, and the trap \texttt{trap\_store\_wire}. We disclose this reduction explicitly:
the three remaining designed cells (\texttt{api\_error\_type}, \texttt{api\_memo\_required},
\texttt{dec\_amount}) are not part of the certified battery. The four-cell selection was made to
retain at least one $W_1$-decisive, one $W_2$-decisive, and the trap cell while keeping the
cross-model spend tractable; we report below that this is best understood as a fixed adversarial
four-cell battery, not a sampled task distribution, and we state the effective support of the
deficiency surrogate accordingly (\S\ref{sec:results-incomp}).

\paragraph{The relevance trap.}
The trap \texttt{trap\_store\_wire} is the confound-proof form of Theorem~\ref{thm:incomparable}. Its
prompt is lexically saturated with encoding vocabulary, so a \emph{query-conditioned} lexical
relevance score ranks $W_2 > W_1$; but its graded requirement is the $W_1$ call contract (it must
post an already-encoded wire string through \texttt{post}), so realized PASS ranks $W_1 > W_2$. This
defeats the query-conditioned relevance score practitioners actually use, not merely a
query-independent ranking. We measure relevance as bag-of-words cosine $\varphi_{\mathrm{lex}}(W,T)$
between a source and the task prompt, computed \emph{before} the model is run. This is the
lexical-relevance instance of Proposition~\ref{prop:trap-constrained}; the empirical mis-rank
is shown for this one concrete scalar (\S\ref{sec:limitations}), while
the theorem (Theorem~\ref{thm:incomparable}) is the general claim for query-independent scalars.

\subsection{Methodology}
\label{sec:methodology}

\paragraph{Models and vendors.}
We evaluate six models across four vendors: Anthropic Haiku-4.5, Sonnet-4.6, and Opus-4.8; OpenAI
GPT-5.5; DeepSeek-V4-Pro; and Moonshot Kimi-K2.6.

\paragraph{Harness and arms.}
The harness (\texttt{tokenbench/measure\_blackwell.py}) is Python-standard-library only (zero
external dependencies); it computes $\dhat_{\Dcal}$ in both directions, the uniform
Clopper--Pearson/Bonferroni certificate, the interference certificate, and the lexical relevance
score, and it implements the launchers for the Anthropic CLI, the Azure Responses API, and
OpenAI-compatible chat completions. Each cell is run under the arms $\{\texttt{none}, W_1, W_2,
W_{1+}\}$, with $W_{1+}^{\mathrm{rev}}$ added for the order-control run. PASS is measured over $n$
completions per cell.

\paragraph{Agentic vs.\ single-shot harness (a stated confound).}
The Anthropic models run \emph{agentically} (\texttt{claude -p}, up to six turns with intermediate
verifier feedback); the non-Anthropic models run as \emph{single-shot} completion plus code
extraction through an OpenAI-style API. Vendor is therefore confounded with harness in this study:
every Anthropic model is agentic and every non-Anthropic model is single-shot. We treat this as a
genuine limitation of the cross-vendor \emph{magnitude} comparisons (\S\ref{sec:limitations}); the
per-model \emph{binary} claims (incomparability certified, relevance mis-rank, anti-monotonicity
certified) are within-model and unaffected. We note that the single-shot setting is arguably the
cleaner probe of whether a model uses supplied context, since it removes the agentic-exploration
pathway.

\paragraph{Certificate parameters.}
All certificates use confidence budget $\eta = 0.10$. Incomparability uses exact Clopper--Pearson
intervals at per-cell level $\alpha = \eta/(2|\Dcal|)$ over the $2|\Dcal|$ proportions
$\{p_{W_1,T}, p_{W_2,T}\}_{T\in \Dcal}$, so by the union bound all hold jointly with probability
$\ge 1-\eta$ and both directions read off the same intervals. The interference certificate uses
$\alpha = \eta/4$ and the threshold $\tau = 0.30$, certifying $\Psi_T \le -\tau$ via the upper bound
of \eqref{eq:psiupper}. We report each certificate as marginally controlled at $\eta$; we do not
claim joint family-wise control across the incomparability, dominance, and interference claims drawn
from a single run (\S\ref{subsec:fwer}, \S\ref{sec:limitations}). The intervals are degenerate at
$p=0$ and $p=1$, which is why the stark $100\%/0\%$ cells certify incomparability at small $n$ while
the interference contrast, which mixes intermediate baselines, does not.

\paragraph{Sampling regime.}
The certificate models each cell as $n$ Bernoulli PASS draws (Assumption~\ref{ass:iid}); the four
arms within a cell are separate completions, hence independent. We flag the exact decoding parameters
(temperature, top-$p$, seeding) and a defense of per-draw independence for the multi-turn agentic arm
as a required protocol addition before submission (\S\ref{sec:limitations}); the binomial model
underlying every $L_{\Dcal}$, $U_{\Dcal}$, and $\Psi_T$ bound depends on it.

\paragraph{Spend.}
The certified record reflects roughly \$170 of Anthropic spend (the live \texttt{claude -p} runs)
plus Azure quota for the non-Anthropic models; the Azure API returns no cost field, so the \$170
figure should not be read as the total cost across all six models.

\subsection{Results}
\label{sec:results}

We report four findings: incomparability (\S\ref{sec:results-incomp}), the relevance mis-rank
(\S\ref{sec:results-misrank}), anti-monotonicity (\S\ref{sec:results-antimono}), and the
model-relative value of private context (\S\ref{sec:results-moat}). Table~\ref{tab:consolidated}
consolidates the per-model verdicts; Tables~\ref{tab:LD}, \ref{tab:interference}, and
\ref{tab:percell} give the underlying quantities.

\begin{table}[t]
\centering
\caption{Consolidated 6-model / 4-vendor verdict. Anthropic rows are agentic (\texttt{claude -p}, up
to six turns); the others are single-shot (\S\ref{sec:methodology}). ``Anti-mono.\ certified'' counts
cells on which the interference upper bound clears $-\tau = -0.30$ at $\eta=0.10$. ``Argorder
\texttt{none}'' is the no-context baseline PASS on \texttt{api\_argorder}, the model-relative-value
spectrum. All figures are real measured runs; no modeled numbers.}
\label{tab:consolidated}
\begin{tabular}{@{}lcccccc@{}}
\toprule
Claim & Haiku & Sonnet & Opus & GPT-5.5 & DeepSeek & Kimi-K2.6 \\
\midrule
Vendor & Anthropic & Anthropic & Anthropic & OpenAI & DeepSeek & Moonshot \\
Harness & agentic & agentic & agentic & single-shot & single-shot & single-shot \\
$n$ per cell & 80 & 40 & 20 & 40 & 40 & 40 \\
Incomp.\ certified ($\ge 90\%$) & \checkmark & \checkmark & \checkmark & \checkmark & \checkmark & \checkmark \\
$L_{\Dcal}(W_1,W_2)$ & $+0.88$ & $+0.76$ & $+0.55$ & $+0.76$ & $+0.347$ & $+0.675$ \\
$L_{\Dcal}(W_2,W_1)$ & $+0.83$ & $+0.76$ & $+0.55$ & $+0.76$ & $+0.715$ & $+0.762$ \\
Relevance mis-rank (trap) & \checkmark & \checkmark & \checkmark & \checkmark & \checkmark & \checkmark \\
Anti-mono.\ certified (\# cells) & 2 & 2 & $0$ (repro.) & 1 & 1 & 1 \\
\texttt{api\_argorder} \texttt{none} & ${\sim}0\%$ & $0\%$ & $15\%$ & $75\%$ & $2\%$ & $45\%$ \\
\bottomrule
\end{tabular}
\end{table}

\subsubsection{Incomparability certified on all six models}
\label{sec:results-incomp}

On every model the empirical deficiency is mutually positive, $\dhat_{\Dcal}(W_1,W_2) > 0$ and
$\dhat_{\Dcal}(W_2,W_1) > 0$, and the uniform lower bounds $L_{\Dcal}$ in both directions clear zero,
certifying empirical incomparability at $\ge 90\%$ (Table~\ref{tab:LD}). The crossover is stark:
$W_1$ solves the API and trap cells and fails encoding, while $W_2$ solves encoding and fails the API
and trap cells. For Haiku at $n=80$ the carrying cells are \texttt{api\_post\_ok} (\texttt{none} $0\%$
/ $W_1$ $98\%$ / $W_2$ $0\%$) and \texttt{enc\_amount} (\texttt{none} $0\%$ / $W_1$ $0\%$ / $W_2$
$100\%$); the certificate strengthened from a marginal $L_{\Dcal} = +0.138$ at $n=10$ to
$L_{\Dcal} = +0.877 / +0.831$ at $n=80$, far clear of zero.

Because incomparability is a two-directional claim, each row of Table~\ref{tab:LD} requires both
bounds positive; the smaller of the two is the binding one. The two carrying cells are one API cell
and one encoding cell, so the effective support of the surrogate is two cells. We therefore describe
$\Dcal$ as a fixed adversarial four-cell battery rather than a sampled task distribution: the
certificate is valid for what it covers, and what it covers is a single hand-built API-versus-encoding
opposition replicated across six models.

\begin{table}[t]
\centering
\caption{Certified incomparability: uniform $(1-\eta)$ lower bounds $L_{\Dcal}$ in both directions
($\eta=0.10$). $L_{\Dcal}(W_1,W_2) > 0$ and $L_{\Dcal}(W_2,W_1) > 0$ together certify empirical
incomparability. Both bounds are read off the same $2|\Dcal|$ Clopper--Pearson intervals
(Proposition~\ref{prop:certificate}).}
\label{tab:LD}
\begin{tabular}{@{}llccc@{}}
\toprule
Model & Vendor & $n$ & $L_{\Dcal}(W_1,W_2)$ & $L_{\Dcal}(W_2,W_1)$ \\
\midrule
Haiku-4.5       & Anthropic & 80 & $+0.877$ & $+0.831$ \\
Sonnet-4.6      & Anthropic & 40 & $+0.762$ & $+0.762$ \\
Opus-4.8        & Anthropic & 20 & $+0.552$ & $+0.552$ \\
GPT-5.5         & OpenAI    & 40 & $+0.762$ & $+0.762$ \\
DeepSeek-V4-Pro & DeepSeek  & 40 & $+0.347$ & $+0.715$ \\
Kimi-K2.6       & Moonshot  & 40 & $+0.675$ & $+0.762$ \\
\bottomrule
\end{tabular}
\end{table}

\paragraph{Why $n=20$ certifies incomparability for Opus but not interference.}
Opus is run at $n=20$, which certifies incomparability ($L_{\Dcal} = +0.552$ both directions) but not
anti-monotonicity. This asymmetry is principled, not an inconsistent threshold (\S\ref{subsec:power}).
The incomparability bound is carried by stark cells whose proportions sit at $0$ or $1$, where the
Clopper--Pearson interval is degenerate and a small $n$ suffices. The interference bound, by contrast,
mixes the $W_{1+}$, $W_1$, $W_2$, and \texttt{none} arms, including the intermediate Opus baselines
(\texttt{api\_argorder} and trap \texttt{none} $= 15\%$), so the four-term Clopper--Pearson sum stays
wide at $n=20$ (upper bounds $-0.08 / -0.11$, not below $-\tau$). Raising Opus to $n=80$ is the one
remaining gap in the consolidated record.

\subsubsection{The relevance score mis-ranks the sources}
\label{sec:results-misrank}

On the trap, query-conditioned lexical relevance ranks $\varphi_{\mathrm{lex}}(W_1) = 0.36 <
\varphi_{\mathrm{lex}}(W_2) = 0.44$, so it prefers $W_2$; yet realized PASS ranks $W_1 \gg W_2$ on all
six models (e.g.\ Haiku $96\%$ vs.\ $0\%$ at $n=80$; Sonnet, GPT-5.5, and Kimi $100\%$ vs.\ $0\%$;
DeepSeek $98\%$ vs.\ $0\%$; Opus $100\%$ vs.\ $10\%$). On the honest cells relevance \emph{agrees}
with PASS, on the API cells relevance ranks $W_1 > W_2$ and so does PASS, and on the encoding cell
relevance ranks $W_2 > W_1$ and so does PASS, so the score fails \emph{only} where utility diverges
from topicality, exactly the regime Theorem~\ref{thm:incomparable} and
Proposition~\ref{prop:trap-constrained} isolate. The $\texttt{none} = 0\%$ floor on the trap confirms
genuine unguessability and $W_2 = 0\%$ confirms that encoding knowledge is useless there.

\subsubsection{Anti-monotonicity: a certified superset that hurts}
\label{sec:results-antimono}

The superset $W_{1+}$, which literally \emph{contains} $W_1$'s signal, sharply degrades the API cells
relative to $W_1$ alone. The per-task interference $\Psi_T = \PASS(W_{1+},T) - \PASS(W_1,T) -
\PASS(W_2,T) + \PASS(\texttt{none},T)$ is certified negative ($\Psi_T \le -\tau$, $\tau=0.30$) on at
least one cell for five of the six models (Table~\ref{tab:interference}). Opus reproduces the effect
with strong point estimates ($\Psi_T = -0.75 / -0.80$) but does not certify at $n=20$. This is a
certified empirical falsification of Blackwell superset-dominance on live models
(Theorem~\ref{thm:antimono}): strictly-more-informative context can hurt, and the deficiency
framework detects it precisely where every monotone scalar must miss it.

\begin{table}[t]
\centering
\caption{Anti-monotonicity (interference) certificate on the carrying cell per model. $\Psi_T =
\PASS(W_{1+}) - \PASS(W_1) - \PASS(W_2) + \PASS(\texttt{none})$; ``Upper'' is the $\eta/4$
Clopper--Pearson upper bound $\overline\Psi_T$ of \eqref{eq:psiupper}; the cell is certified harmful
when Upper $\le -0.30$. The $W_1 \to W_{1+}$ column reports the PASS collapse. The certificate clears
only on cells whose \texttt{none} baseline floors near zero (the leaky-baseline limitation;
Remark~\ref{rem:leaky}).}
\label{tab:interference}
\begin{tabular}{@{}llcccc@{}}
\toprule
Model & Carrying cell & $W_1 \to W_{1+}$ & $\Psi_T$ & Upper & Certified \\
\midrule
Haiku-4.5       & \texttt{api\_post\_ok}            & $98\% \to 24\%$  & $-0.74$ & $-0.49$ & \checkmark \\
Haiku-4.5       & \texttt{api\_argorder} (hardened) & $100\% \to 1\%$  & $-0.97$ & $-0.79$ & \checkmark \\
Sonnet-4.6      & \texttt{api\_post\_ok}            & $100\% \to 0\%$  & $-1.00$ & $-0.69$ & \checkmark \\
Sonnet-4.6      & \texttt{api\_argorder}            & $100\% \to 0\%$  & $-1.00$ & $-0.69$ & \checkmark \\
GPT-5.5         & \texttt{api\_post\_ok}            & $100\% \to 0\%$  & $-0.95$ & $-0.60$ & \checkmark \\
DeepSeek-V4-Pro & \texttt{api\_post\_ok}            & $75\% \to 0\%$   & $-0.75$ & $-0.36$ & \checkmark \\
Kimi-K2.6       & \texttt{api\_post\_ok}            & $100\% \to 0\%$  & $-1.00$ & $-0.69$ & \checkmark \\
Opus-4.8        & \texttt{api\_post\_ok}            & $100\% \to 25\%$ & $-0.75$ & $-0.08$ &, ($n=20$) \\
Opus-4.8        & \texttt{api\_argorder}            & $100\% \to 5\%$  & $-0.80$ & $-0.11$ & no ($n=20$) \\
\bottomrule
\end{tabular}
\end{table}

\paragraph{The leaky-baseline limitation is a stated protocol.}
The interference upper bound includes the term $+\mathrm{CP}_{\mathrm{hi}}(p_{\texttt{none}})$, so it
can certify harm only on cells whose no-context baseline floors near zero. This is exactly why
anti-monotonicity certifies on \texttt{api\_argorder} (hardened) for Haiku (\texttt{none} dropped from
$40\%$ to $1\%$ after hardening, $W_1$ still $100\%$) but not on \texttt{api\_argorder} for GPT-5.5
(\texttt{none} $= 75\%$), Kimi ($45\%$), or Opus ($15\%$). We state this as a protocol up front
(Remark~\ref{rem:leaky}): we certify anti-monotonicity only on cells whose \texttt{none} baseline
floors near zero, because a leaky baseline mechanically shrinks $\Psi_T$'s certifiable magnitude.
Table~\ref{tab:interference} therefore reports, per model, which cell carried the certificate and the
collapse on it. We also report the $W_1$ ceiling alongside every $\Psi_T$: DeepSeek's $W_1$ arm
reaches only $75\%$ on \texttt{api\_post\_ok} (an honest capability difference, vs.\ ${\sim}100\%$ for
Claude and GPT), so its $\Psi_T = -0.75$ is computed from a lower starting point and is not directly
magnitude-comparable to Sonnet's $100\% \to 0\%$.

\paragraph{Order control.}
Reordering the superset does not restore performance. At $n=80$ on Haiku, \texttt{api\_post\_ok} gives
$W_{1+}$ $24\%$ and $W_{1+}^{\mathrm{rev}}$ $1\%$ (vs.\ $W_1$ $98\%$), and \texttt{api\_argorder} gives
$W_{1+}$ $24\%$ and $W_{1+}^{\mathrm{rev}}$ $44\%$ (vs.\ $W_1$ $100\%$). On one cell the reversed order
is worse and on the other it is better, so order does have a directional effect; but \emph{neither}
ordering recovers $W_1$ alone. The defensible reading is that reordering does not cure the harm, hence
the collapse is not a lost-in-the-middle / recency artifact, not that order is irrelevant.

\paragraph{Runtime fingerprint (a descriptive correlate).}
The harmful $W_{1+}$ arm consistently emits more output tokens than $W_1$: $+106\% / +87\% / +252\%$
across the Haiku collapse and encoding cells (and $+12\%$ even on the still-passing encoding superset
cell), Sonnet $+120\% / +131\%$, Opus $+204\% / +204\%$, GPT-5.5 $+281\% / +274\%$, DeepSeek $+93\% /
+166\%$, Kimi $+203\% / +249\%$. We report this as a consistent descriptive correlate, not a
mechanism: more output on the failing arm could be effect rather than cause, and the $+12\%$ inflation
on a cell that still passes undercuts a strict ``verbose thrashing causes the collapse'' reading.

\paragraph{Robust, but not monotone in capability.}
The harm is large and robust across the Anthropic capability ladder but is \emph{not} ordered by
capability. On the two API cells the $W_1 \to W_{1+}$ collapse is Haiku $100\% \to 24\% / 1\%$, Sonnet
$100\% \to 0\% / 0\%$, and Opus $100\% \to 25\% / 5\%$: the middle model (Sonnet) shows the starkest
collapse and the strongest model (Opus) less so. An earlier two-point reading suggested the harm
``intensifies with strength''; the third point falsifies that, and we explicitly do not claim a
capability trend, only robustness across the ladder.

\subsubsection{The model-relative value of private context}
\label{sec:results-moat}

The ``private'' conventions are genuinely unguessable for some vendors and partially known to others,
and this varies continuously. The \texttt{none}-arm baseline on \texttt{api\_argorder} forms a clean
spectrum: Claude ${\sim}0\%$, DeepSeek $2\%$, Opus $15\%$, Kimi $45\%$, GPT-5.5 $75\%$
(Table~\ref{tab:consolidated}). Thus the value of the private context, the model-relative quantity
$I(S;W\mid\theta)$ of \S\ref{subsec:theta}, varies continuously across models. It tracks a vendor's
training rather than raw capability: DeepSeek, a different vendor from Anthropic, nonetheless has a
tight prior (baselines $0/2/0/0\%$ across \texttt{api\_post\_ok}/\texttt{api\_argorder}/\texttt{enc\_amount}/trap)
like Claude, while GPT-5.5's prior already holds much of the convention (\texttt{api\_argorder} $75\%$,
\texttt{enc\_amount} $55\%$). Sonnet's baselines all floor at $0\%$, so even the stronger Anthropic
model does not guess the convention; the value of private context survives the capability ladder. The
certified per-model claims rest on cells unguessable \emph{for the model under test} (e.g.\ GPT-5.5
\texttt{api\_post\_ok} \texttt{none} $= 5\%$, trap $= 0\%$). We note that this cross-vendor magnitude
comparison is the part most exposed to the agentic-vs-single-shot confound (\S\ref{sec:limitations}):
vendor and harness are not separated in the present runs.

\paragraph{Auditable per-cell record.}
Table~\ref{tab:percell} gives the per-arm PASS rates for the carrying cells so each certificate can be
reconstructed; the full per-cell, per-model record is in the supplementary material and the run logs
(\texttt{blackwell\_*\_n*.\{log,json\}}).

\begin{table}[t]
\centering
\caption{Per-cell PASS rates ($\%$) for the carrying cells, by arm. Anthropic rows agentic, others
single-shot. ``, '' marks an arm not in that run. Hardened \texttt{api\_argorder} (Haiku) lowered the
\texttt{none} baseline from $40\%$ to $1\%$ while leaving $W_1$ at $100\%$.}
\label{tab:percell}
\begin{tabular}{@{}llccccc@{}}
\toprule
Model & Cell & \texttt{none} & $W_1$ & $W_2$ & $W_{1+}$ & $W_{1+}^{\mathrm{rev}}$ \\
\midrule
Haiku-4.5       & \texttt{api\_post\_ok}           & 0  & 98  & 0   & 24  & 1   \\
Haiku-4.5       & \texttt{api\_argorder} (hard.)   & 1  & 100 & 0   & 1   & 44  \\
Haiku-4.5       & \texttt{enc\_amount}             & 0  & 0   & 100 & -- & -- \\
Haiku-4.5       & \texttt{trap\_store\_wire}       & 0  & 96  & 0   & -- & -- \\
Sonnet-4.6      & \texttt{api\_post\_ok}           & 0  & 100 & 0   & 0   & -- \\
Sonnet-4.6      & \texttt{api\_argorder}           & 0  & 100 & 0   & 0   & -- \\
Sonnet-4.6      & \texttt{trap\_store\_wire}       & 0  & 100 & 0   & -- & -- \\
Opus-4.8        & \texttt{api\_post\_ok}           & 0  & 100 & 0   & 25  & -- \\
Opus-4.8        & \texttt{api\_argorder}           & 15 & 100 & 0   & 5   & -- \\
Opus-4.8        & \texttt{trap\_store\_wire}       & 15 & 100 & 10  & -- & -- \\
GPT-5.5         & \texttt{api\_post\_ok}           & 5  & 100 & 0   & 0   & -- \\
GPT-5.5         & \texttt{api\_argorder}           & 75 & -- & -- & -- & -- \\
GPT-5.5         & \texttt{enc\_amount}             & 55 & -- & -- & -- & -- \\
GPT-5.5         & \texttt{trap\_store\_wire}       & 0  & 100 & 0   & -- & -- \\
DeepSeek-V4-Pro & \texttt{api\_post\_ok}           & 0  & 75  & 0   & 0   & -- \\
DeepSeek-V4-Pro & \texttt{api\_argorder}           & 2  & 68  & 0   & -- & -- \\
DeepSeek-V4-Pro & \texttt{trap\_store\_wire}       & 0  & 98  & 0   & -- & -- \\
Kimi-K2.6       & \texttt{api\_post\_ok}           & 0  & 100 & 0   & 0   & -- \\
Kimi-K2.6       & \texttt{api\_argorder}           & 45 & -- & -- & -- & -- \\
Kimi-K2.6       & \texttt{trap\_store\_wire}       & 0  & 100 & 0   & -- & -- \\
\bottomrule
\end{tabular}
\end{table}

\subsection{Generality (a second pair) and harness robustness (single-shot)}
\label{sec:results-generality}

Two follow-up experiments address the two largest threats to validity directly.

\paragraph{A second incomparable pair, different domain.}
We repeated the protocol with a second hand-built pair on a private \texttt{cache} module: $V_1$, the
\texttt{cache.put(key, value, *, ttl)} call contract (argument order, required keyword-only
\texttt{ttl}, a True/False ``newly inserted'' flag, a \texttt{CacheError} type), versus $V_2$, a
key-normalization invariant (a fixed \texttt{kx7-} namespace prefix and an underscore rule). On Haiku
at $n=40$ the pair is again certified incomparable ($L_{\Dcal} = +0.774 / +0.496$, both directions clear
of zero) with a clean double crossover ($V_1$ solves the contract and trap cells, $V_2$ solves
key-normalization, all none-baselines floor at $0\%$), and the dominance control again certifies
one-directionally. Anti-monotonicity reproduces in direction (the superset collapses the contract cell
$80\%\!\to\!28\%$, $\Psi_T=-0.53$) but does not certify at $n=40$ (the $V_1$ rescue is only $80\%$, a
smaller gap than pair~1; it would require $n\!\approx\!80$). The relevance trap for this pair was a
null, its prompt vocabulary leaned toward $V_1$ ($\varphi(V_1)=0.39 > \varphi(V_2)=0.25$), so
relevance \emph{agreed} with PASS, so the strong mis-rank remains a pair-1 result. The contribution of
the second pair is therefore \emph{generality of the incomparability and dominance results to a new
domain}, not a second mis-rank.

\paragraph{Single-shot breaks the harness confound.}
To separate vendor from harness, we ran an Anthropic model (Haiku) in \emph{single-shot completion}
mode, one turn, no tools, code extracted from the returned text, the same regime as the non-Anthropic
models. At $n=40$ the core results survive the non-agentic harness: incomparability is certified
($L_{\Dcal} = +0.762$ both directions), the relevance mis-rank reproduces (trap PASS $W_1$ $90\%$ vs
$W_2$ $0\%$), the dominance control certifies, and anti-monotonicity is certified on
\texttt{api\_argorder} ($\Psi_T = -0.93$, upper bound $-0.56$). The one regime-dependent quantity is
absolute execution reliability: single-shot $W_1$ rescue on \texttt{api\_post\_ok} is $52\%$ versus
$98\%$ agentic (no retries to self-correct), so that cell's interference does not certify single-shot
while its order-theoretic claims do. The order-theoretic findings, certified incomparability, relevance
mis-rank, dominance, and at least one certified interference cell, are therefore \emph{not} artifacts of
the agentic harness.

% ==========================================================================
\section{Discussion}
\label{sec:discussion}
% ==========================================================================

Our results recast a piece of practitioner folklore, ``relevance does not transfer across
tasks'', as a structural fact about context sources. When two sources are Blackwell-incomparable, the
order over them is genuinely partial, and any scalar score, by imposing a total order, must mis-rank
them on some task. The empirical program does not merely illustrate this; it instantiates a single
hand-built incomparable pair (a private API-call contract $W_1$ versus a private wire-encoding
invariant $W_2$ on the same \texttt{ledger} module) and shows that the resulting double crossover, the
relevance mis-rank, and an anti-monotonicity collapse all reproduce across six models and four
vendors. We discuss what the evidence does and does not license, and where the
formalization is a model rather than a proof of the realized object.

\subsection{Theorems are about the idealized decision-maker; the experiments measure a fixed rule}
\label{sec:disc-bridge}

The cleanest way to read the theory is as a two-layer claim, and the two layers must not be conflated.
Corollary~\ref{cor:c1} (universal context dominance) and Theorem~\ref{thm:incomparable} (no scalar
ranks incomparable sources) are statements about the \emph{Bayes-optimal} user of a source: ``$W_1
\Bord W_2$'' means the optimal achievable value under $W_1$ dominates that under $W_2$ for every
prior, bounded utility, and decision rule~\citep{blackwell1953}. The operational quantity in every
experiment, by contrast, is $\PASS(W,T)=\mathbb{E}[v_T(\delta_M(T,w))]$ for a single \emph{fixed,
sub-optimal} decision rule $\delta_M$, the language model. ``Optimal-user-of-$W_2$ beats
optimal-user-of-$W_1$ on some decision problem'' does not, by itself, imply ``the fixed model's PASS
under $W_2$ exceeds its PASS under $W_1$ on task $T_b$.'' Lemma~\ref{lem:bridge} supplies the bridge
in the direction we need (certified $(M,\Dcal)$-incomparability implies Blackwell incomparability) and
isolates the value-attainment hypothesis (Assumption~\ref{ass:attain}) required for the converse. We
therefore read the contribution as a $\delta_M$-restricted, $\Dcal$-restricted incomparability
statement, certified directly from PASS data, that is \emph{consistent with} and motivated by the
idealized Blackwell incomparability of the underlying coordinates, and, via the bridge, earns the
structural conclusion rather than assuming it.

\subsection{The estimator is a regret-gap surrogate, not a Le~Cam deficiency}
\label{sec:disc-deficiency}

Genuine Le~Cam deficiency $\delta(\mathcal{E},\mathcal{F})$ minimizes over garblings and takes a
supremum over a class of loss functions~\citep{lecam1986}. Our surrogate $\dhat_{\Dcal}$
\eqref{eq:estimator} performs \emph{no} garbling minimization, fixes a single sub-optimal rule
$\delta_M$, and ranges over a finite task battery $\Dcal$. It is a one-sided decision-restricted
\emph{value/regret gap}, not a deficiency in Le~Cam's sense, and we do not prove a two-sided
inequality relating it to the true deficiency (for a fixed sub-optimal rule none holds in general in
either direction; Remark~\ref{rem:notlecam}). We accordingly read $\dhat_{\Dcal}$ as a
``decision-restricted value-deficiency surrogate'' and demote the Le~Cam billing to an analogy that
motivates the construction. What is rigorous is the \emph{certificate}: the distribution-free uniform
Clopper--Pearson~\citep{clopperpearson1934} plus Bonferroni~\citep{bonferroni1936} lower bound
$L_{\Dcal}$ on $\sup_T[\PASS(W_2,T)-\PASS(W_1,T)]$, which we regard as the strongest methodological
component and which survives independently of the deficiency framing.

\subsection{What the certificate covers, and what it does not}
\label{sec:disc-certificate}

The uniform certificate is the part of the paper we are most confident in. Its correctness rests on
paying the multiplicity of the $\sup_T$ selection via a union bound over $2|\Dcal|$ exact intervals
rather than absorbing it into a single-sample sup inequality; a naive DKW bound under-covers precisely
because the object is a maximum of \emph{differences of two binomial means across distinct cells}, not
the sup-deviation of one empirical CDF from its own (Remark~\ref{rem:dkw}). Two qualifications
follow. First, the effective support carrying the certificate is small: although $\Dcal$ has four
cells, $L_{\Dcal}(W_1,W_2)$ is realized on one encoding cell and $L_{\Dcal}(W_2,W_1)$ on one API cell,
so the certified incomparability is a two-cell crossover replicated across models, not a sample from a
task distribution. We therefore describe $\Dcal$ as a fixed adversarial battery rather than a ``task
distribution.'' Second, each certified claim per model is controlled \emph{marginally}, not jointly:
incomparability (at $\alpha=\eta/(2|\Dcal|)$ over $2|\Dcal|$ intervals), the dominance control, and the
interference certificate (at $\alpha=\eta/4$) draw from the same run at the same $\eta=0.10$, so the
family-wise error across all certified claims on a model is not simultaneously controlled
(\S\ref{subsec:fwer}). Readers should treat each marginal certificate as valid in isolation; a strict
joint statement would require a wider global Bonferroni budget.

\subsection{Anti-monotonicity is real, certified, but not capability-ordered}
\label{sec:disc-antimono}

The interference object $\Psi_T = \PASS(W_{1+},T) - \PASS(W_1,T) - \PASS(W_2,T) + \PASS(\text{none},T)$
measures whether appending the (irrelevant-to-this-task) $W_2$ to a sufficient $W_1$ harms
performance. We certify $\Psi_T \le -0.30$ on at least one cell for five of six models
(Table~\ref{tab:interference}), with collapses as stark as $100\% \to 0\%$ ($\Psi_T=-1.00$, upper
bound $-0.69$) on Sonnet and Kimi. This is a certified empirical falsification of
superset-dominance for these fixed rules (Theorem~\ref{thm:antimono}): a literal superset of a
sufficient signal can sharply reduce PASS, which no monotone-submodular scalar context-value score can
represent.

Two corrections follow from the data. \textbf{(i) The harm is not monotone in
capability.} An early two-point read (Haiku then Sonnet) suggested the collapse ``intensifies with
model strength''; adding the third Anthropic model (Opus) falsified this. On \texttt{api\_post\_ok}
the collapse is $100\to24\%$ (Haiku), $100\to0\%$ (Sonnet), $100\to25\%$ (Opus), and on
\texttt{api\_argorder} $100\to1\%$, $100\to0\%$, $100\to5\%$ respectively, the middle model (Sonnet)
is the starkest and the strongest (Opus) is less so. The defensible claim is robustness across the
capability ladder, \emph{not} a capability trend.
\textbf{(ii) The interference certificate is structurally cell-selected.} Its upper bound includes
$+\mathrm{CP}_{\mathrm{hi}}(p_{\text{none}})$, so it can only certify harm on cells whose none-baseline
floors near zero. This is exactly the set of cells that are unguessable for the model under test, which
couples the certifiable-harm set to the model-relative value of private context (\S\ref{sec:disc-moat}).
We treat this as a protocol (Remark~\ref{rem:leaky}): we certify anti-monotonicity only on
floor-baseline cells, because a leaky baseline
mechanically shrinks the certifiable $\Psi_T$. On \texttt{api\_argorder} the harm is certified for
Haiku (hardened baseline $1\%$) and Sonnet ($0\%$) but not for GPT-5.5 (baseline $75\%$), DeepSeek,
Kimi ($45\%$), or Opus ($15\%$).

\subsection{The value of private context is model-relative and tracks vendor training}
\label{sec:disc-moat}

The value of a private convention, $I(S;W\mid\theta)$, is itself a model-relative ($\theta$-dependent)
object, and the data give a clean continuous map of it. The none-arm baseline on \texttt{api\_argorder}
, how often a model reproduces the private convention with no context supplied, forms a spectrum:
Claude $\sim\!0\%$, DeepSeek $2\%$, Opus $15\%$, Kimi $45\%$, GPT-5.5 $75\%$. We do not claim this
tracks capability. The decisive piece of evidence is DeepSeek: it has a tight, near-zero prior like
Claude despite being a different vendor, while GPT-5.5's prior is wide ($55\%$--$75\%$ guessing on
\texttt{api\_argorder}/\texttt{enc\_amount}). A ``stronger-model-knows-more'' account predicts the
spectrum to follow capability; it does not. We read $I(S;W\mid\theta)$ as tracking a vendor's training
distribution rather than raw capability, which makes the theory's $\theta$-dependence concrete and
caps which cells certify per model.

\subsection{Mechanism: the runtime fingerprint is a correlate, not an explanation}
\label{sec:disc-mechanism}

Every collapse is accompanied by output-token inflation on the harmful arm: $W_{1+}$ emits between
$+87\%$ and $+252\%$ more output tokens than $W_1$ on the API collapse cells (e.g.\ Haiku
$+106\%/+87\%/+252\%$; Sonnet $+120\%/+131\%$; GPT-5.5 $+281\%/+274\%$). One might read this as
``verbose thrashing'' causing the failure, but the relationship is correlational and uncontrolled: the
same superset arm also inflates output by $+12\%$ on the \emph{still-passing} encoding cell, which
undercuts a clean ``verbosity causes collapse'' account. We therefore treat the token fingerprint as a
reproducible correlate of the harmful arm, not a mechanism.

% ==========================================================================
\section{Limitations and threats to validity}
\label{sec:limitations}
% ==========================================================================

We collect the scope restrictions here.

\paragraph{Scaling the selector is the central open problem.}
\textsc{Probe-Select} (\S\ref{sec:method}) is curation-scale: building the certified order on $m$
candidate sources costs $m\!\cdot\!|\Dcal|\!\cdot\!n$ verifier-gated model calls, which is practical
for choosing among a handful of private conventions or curating a few-shot pool, but not for
retrieval over a corpus of $10^6$ passages. We prove the selector's properties and validate the
estimator that powers it (\S\ref{sec:experiments}), but we do \emph{not} benchmark a corpus-scale
system against strong dense rerankers. Closing this gap, most plausibly by amortizing the probe
into a learned deficiency predictor, or by using \textsc{Probe-Select} as a partial-order re-ranking
head over a cheap scalar shortlist, is, in our view, the most important direction this work opens,
and the one on which its practical impact depends.

\paragraph{Source-pair generality (partially addressed).}
The cross-model record (six models) rests on one adversarially hand-built pair, the API contract $W_1$
versus the wire-encoding invariant $W_2$. To test whether the partial order is a \emph{pattern} rather
than a one-off, we ran a second pair in a different domain, a \texttt{cache}-module call contract $V_1$
versus a key-normalization invariant $V_2$, which is again certified incomparable on Haiku
($L_{\Dcal}=+0.774/+0.496$) with a certified one-directional dominance control
(\S\ref{sec:results-generality}). Incomparability and dominance therefore hold for \emph{two} pairs in
\emph{two} domains, not one. The remaining content gaps are narrower: pair~2's anti-monotonicity is only
directional at $n=40$ (would certify at $n\!\approx\!80$), pair~2's relevance trap was a null (its prompt
was not lexically $V_2$-leaning), and both effects on pair~2 are shown for a single model (Haiku), so the
\emph{cross-vendor} record is still single-pair. A third domain and a vendor-general strong mis-rank
remain inexpensive next steps.

\paragraph{Vendor is perfectly confounded with harness.}
The Anthropic models (Haiku, Sonnet, Opus) run agentically via \texttt{claude -p} with up to six turns
and intermediate verifier feedback; the non-Anthropic models (GPT-5.5, DeepSeek, Kimi) run as
single-shot completion followed by code extraction over an OpenAI-style API. Every Anthropic row is
therefore agentic and every non-Anthropic row single-shot, so ``cross-vendor'' and ``cross-harness''
cannot be separated in the six-model dataset. The binary per-model claims (incomparability certified;
relevance mis-rank; at least one certified anti-monotonicity cell) are within-model and do not depend on
cross-arm comparability. To break the confound for the order-theoretic claims, we ran one Anthropic
model (Haiku) in the single-shot regime (\S\ref{sec:results-generality}): incomparability, the relevance
mis-rank, dominance, and one certified interference cell all survive the non-agentic harness, so these
findings are not artifacts of the agentic arm. The cross-model \emph{magnitude} comparisons, the
$0/2/15/45/75\%$ value spectrum and the $\Psi_T$ severity ladder, still depend on cross-arm
comparability (the harness-crossed run was on one model, and absolute execution reliability \emph{is}
regime-dependent: $W_1$ rescue $98\%$ agentic vs.\ $52\%$ single-shot on \texttt{api\_post\_ok}), so
those magnitude comparisons should still be read as suggestive. We argue single-shot is arguably the
\emph{cleaner} probe of whether a model uses the supplied context (no agentic-exploration confound).

\paragraph{Decision-restricted, $\Dcal$-restricted scope.}
$\dhat_{\Dcal}$ certifies dominance and incomparability \emph{for a given model on a fixed task
battery}, not the full universal Blackwell order. This restriction is exactly what makes the object
computable from PASS data, and within $\Dcal$ it remains a decision-independent-of-task guarantee, but
it is not the universal order, and we do not claim it is.

\paragraph{Sampling regime and per-draw independence.}
The certificate assumes $n$ i.i.d.\ Bernoulli PASS outcomes per cell (Assumption~\ref{ass:iid}). We
report the exact decoding parameters, seeding policy, and turn budget in \S\ref{sec:repro}, and we note
the assumption is most delicate for the agentic arm, where each draw is a multi-turn trajectory; we
treat each completion as a separate draw and depend on per-draw independence for the binomial intervals.
The asymmetry by which $n=20$ certifies incomparability for Opus but not its interference is principled
rather than ad hoc (\S\ref{subsec:power}): stark $0/1$ cells make the Clopper--Pearson endpoints
collapse, so few draws suffice, whereas the interference contrast mixes intermediate baselines and needs
more, so the two certificates have different effective-$n$ requirements by design.

\paragraph{Leaky baselines bound the certifiable set.}
As in \S\ref{sec:disc-antimono}, anti-monotonicity can only be certified on cells whose none-baseline
floors near zero. We treat ``certify only on floor-baseline cells'' as a stated protocol
(Remark~\ref{rem:leaky}), and report per model which cell carried the certificate and its baseline
(Table~\ref{tab:consolidated}, Table~\ref{tab:interference}).

\paragraph{Relevance is demonstrated only for lexical cosine.}
The mis-rank ($\varphi(W_1)=0.36 < \varphi(W_2)=0.44$ on \texttt{trap\_store\_wire}, while realized PASS
is $W_1\!\gg\!W_2$, e.g.\ $96$--$100\%$ vs $0\%$) is shown for bag-of-words cosine relevance. The theorem
rules out any \emph{query-independent} scalar (Theorem~\ref{thm:incomparable}); the trap is an
\emph{empirical} patch against \emph{query-conditioned} relevance (Proposition~\ref{prop:trap-constrained})
and is shown for one concrete scalar. A learned dense reranker might or might not mis-rank the trap;
reproducing the mis-rank with an off-the-shelf embedding reranker would strengthen the strong-form claim,
which we currently restrict to lexical relevance.

\paragraph{Marginal, not joint, certification.}
The incomparability, dominance, and interference certificates drawn from one run are each controlled at
$\eta = 0.10$ in isolation; we do not claim joint family-wise control across all certified claims per
model (\S\ref{subsec:fwer}).

\paragraph{Synergy (positive interference) is untested.}
By construction $W_1$ and $W_2$ are task-disjoint, so the experiments measure negative interference only;
a complementary pair (two facts a task needs \emph{jointly}) would test whether interference is
sign-indefinite (i.e.\ whether a superset can also \emph{help} beyond its parts). We do not measure this
and make no claim about positive synergy.

% ==========================================================================
\section{Conclusion}
\label{sec:conclusion}
% ==========================================================================

Context selection for language models is, structurally, a partial-order problem rather than a scalar
one. When two context sources are Blackwell-incomparable, no scalar score can rank them correctly for
all tasks (Theorem~\ref{thm:incomparable}), which gives a formal account of the field's ``relevance does
not transfer'' folklore. We made the order measurable with a decision-restricted value-deficiency
surrogate $\dhat_{\Dcal}$ equipped with a distribution-free uniform Clopper--Pearson/Bonferroni
certificate, and we instantiated the theory on six models across four vendors. On every model we certify
($\ge 90\%$) that two private context sources are empirically incomparable; on every model a real
query-conditioned lexical relevance score mis-ranks them; and on five of six models we certify a superset
of a sufficient source sharply degrading PASS (a collapse as deep as $100\%\to 0\%$), a certified
falsification of superset-monotonicity for these fixed rules. The value of private context
$I(S;W\mid\theta)$ is itself model-relative, forming a continuous $0\%$--$75\%$ spectrum that tracks
vendor training rather than raw capability.

On scope: the theorems concern the idealized Bayes-optimal user of a source; the
realized crossover, anti-monotonicity, and value-of-private-context are properties of \emph{fixed} models
measured directly, bridged to the structural claim by Lemma~\ref{lem:bridge}. The estimator is a
regret-gap surrogate, not a Le~Cam deficiency; the certificate, which we regard as the most robust
contribution, is correct and survives that re-labeling. The evidence now spans two hand-built source
pairs in two domains and a single-shot replication that preserves the order-theoretic claims; the
cross-vendor record remains single-pair, vendor is confounded with harness for the magnitude
comparisons, and synergy is untested. The practical reading is modest and actionable: selection should
be treated as partial-order-aware rather than scalar top-$k$, the right operational object is a certified
deficiency surrogate rather than a relevance scalar, and ``the best context'' should be expected to be
model-relative and only partially ordered. The remaining strengthenings, a vendor-general strong
mis-rank, an embedding-reranker relevance baseline, a third domain, and certifying pair-2
anti-monotonicity at $n\!\approx\!80$, are the natural next steps, and the experimental record indicates
each is cheap to run.

% ==========================================================================
\section{Reproducibility}
\label{sec:repro}
% ==========================================================================

All reported numbers come from real model calls whose outputs are checked by executable verifiers; there
are no modeled or aspirational figures. The harness (\texttt{tokenbench/measure\_blackwell.py}) is
Python-standard-library only (no third-party dependencies; the Clopper--Pearson endpoints use a stdlib
regularized incomplete beta), is cross-OS, and implements the surrogate $\dhat_{\Dcal}$, the uniform
incomparability certificate (\texttt{certify\_incomparable}, \texttt{clopper\_pearson}), the dominance
control (\texttt{certify\_dominance}), the interference certificate, the sample-size helpers
(\texttt{required\_n}, \texttt{required\_n\_stark}), and the lexical-relevance scalar
(\texttt{lexical\_relevance}). It includes launchers for the Anthropic CLI (\texttt{claude -p}), the
Azure Responses API, and OpenAI-compatible chat completions; the launcher auto-detects API style from the
endpoint, and all API keys are read from the environment (never written to a file).

\paragraph{Protocol.} The verifier stubs a fake \texttt{ledger} module that accepts \emph{only} the exact
private contract (it rejects wrong positional arg order, a missing \texttt{memo=} keyword, and the wrong
exception type); return code $0$ is PASS. The four certified cells are \texttt{api\_post\_ok},
\texttt{api\_argorder} (hardened so \texttt{memo} arrives via a parameter, forcing the keyword-only fact
and dropping its baseline from $40\%$ to $1\%$), \texttt{enc\_amount}, and \texttt{trap\_store\_wire};
arms are $\{\text{none}, W_1, W_2, W_{1+}\}$ plus $W_{1+}^{\text{rev}}=W_1\!+\!\!+\,W_2$ for the order
control. PASS is measured over $n$ i.i.d.\ completions per cell and the certificate uses $\eta=0.10$. The
Anthropic arms run agentically (\texttt{claude -p}, up to six turns); the non-Anthropic arms run
single-shot (one completion + code extraction). We note for transparency that this design
choice, agentic for Anthropic, single-shot for the others, is the vendor/harness confound discussed in
\S\ref{sec:limitations}.

\paragraph{Artifacts.} Per-model run logs and result JSONs are in the repository under the pattern
\texttt{blackwell\_*\_n*.\{log,json\}}, including \texttt{blackwell\_n80\_run.\{log,json\}} (the $n=80$
Haiku bundled certification, $1600$ calls), \texttt{blackwell\_argorder\_n80.\{json,log\}} (the hardened
second anti-monotonicity cell), \texttt{blackwell\_sonnet\_n40}, \texttt{blackwell\_opus\_n20},
\texttt{blackwell\_gpt55\_n40}, \texttt{blackwell\_deepseek\_n40}, and \texttt{blackwell\_kimi\_n40}. The
full experimental record with exact per-cell figures is in \texttt{EXPERIMENT-BLACKWELL.md} (\S\S10--13)
and the theory hardening (formalization and proof sketches) in \texttt{BLACKWELL.md} (\S10).

\paragraph{Compute and cost.} Total measured spend was approximately \$170 on Anthropic
(\texttt{claude -p}): roughly \$123 across the Haiku and Sonnet anchors plus the Opus run (\$49.84) and
the order-control and second-cell runs. The non-Anthropic runs (GPT-5.5, DeepSeek, Kimi) used Azure
quota; the Azure Responses API returns no cost field, so those rows report a \$0 cost artifact while token
counts are captured, the \$170 figure is the Anthropic spend, not the total cost of all six models. The
two highest-value strengthenings flagged in the record (a second source pair on Haiku, and one
harness-crossed run) are each estimated in the tens of dollars.

\paragraph{Reference verification (author note).} \emph{The bibliography was assembled in part with
automated literature search, which has produced incorrect arXiv identifiers and claims in this project
before, so every entry was subsequently web-verified against arXiv/DBLP/publisher records. The outcome:
zero fabricated references, all entries resolve to real papers, but eight had incorrect fields (now
corrected; e.g.\ the BED-LLM identifier was repointed from \texttt{2604.00414}, a different paper, to
\texttt{2508.21184}), and four in-text claims were re-sourced or softened where the cited paper did not
support the sentence (most importantly, the ``relevance does not transfer across tasks'' folklore is now
stated without attaching it to a paper that is not about it). A final line-by-line check against the
primary sources is still recommended before submission; the classics (Blackwell 1953; Le~Cam 1964/1986;
Ethayarajh et al.\ 2022) are author-verified.}

% ==========================================================================
\bibliographystyle{plainnat}
\bibliography{blackwell-paper}
% ==========================================================================

\end{document}
